% BEGIN AUTHOR FILE sec-introduction.tex
\section{Introduction}
\label{sec:introduction}

There are two primary tasks in topological data analysis (TDA)~\cite{DeyWang2022,EdelsbrunnerHarer2010,Zomorodian2005}: reconstruction and inference. 
In a typical TDA pipeline, the data is given as a point cloud in $\Rspace^N$.  
A ``geometric shape'' $K$ in $\Rspace^N$ is reconstructed from the point cloud, usually as a simplicial complex, and $K$ is taken to represent the (unknown) space $X$ from which the data is sampled. 
$X$ is then studied using $K$ as a surrogate for properties that are invariant under invertible mappings (technically, homeomorphisms). 
The deduced ``topological shape'' is not specific to the complex $K$ or the space $X$, but is a feature of the homeomorphism type of $K$ or $X$. 
For example, a standard round circle has the same topological shape as any closed loop such as a knot in the Euclidean space. 
Although topological properties of $X$ alone are not sufficient to reconstruct $X$ exactly, they are among a few global features of $X$ that can be inferred from the data sampled from $X$. 

An (ordinary) persistence barcode~\cite{CarlssonZomorodianCollins2004,Ghrist2008} (or equivalently, a persistence diagram~\cite{ChazalCohenSteinerGlisse2009, Cohen-SteinerEdelsbrunnerHarer2005, EdelsbrunnerLetscherZomorodian2002}) captures the evolution of homological features in a filtration constructed from a simplicial complex $K$. 
It consists of a multi-set of intervals in the extended real line, where the start and end points of an interval (i.e.,~a bar) are the birth and death times of a homological feature in the filtration. 
Equivalently, a persistence diagram is a multi-set of points in the extended plane, where a point in the persistence diagram encodes the birth and death time of a homological feature. 

A main drawback of ordinary persistent homology is that it does not provide canonical choices of geometric representatives for each bar in the barcode. 
Specifically, there are two levels of choices in assigning representatives for bars in a persistence barcode: first, we choose a basis for persistent homology in a consistent way across the filtration; second, we choose a cycle representative inside each homology class in the basis. Both levels of choices involve choosing among significantly distinct geometric features of data to represent the same bar in the barcode. Since such choices are not canonical, it can be challenging to interpret their underlying geometric meaning. In addition, while existing works~\cite{Oleksiy2010,DeyHouMandal2020, Obayashi2018} compute optimal representatives for ordinary persistence, the optimality would rely on a choice of weights for the simplices and on the optimality criterion, leading to different shapes. Moreover, computing the optimal representatives is NP-hard in many interesting situations  ~\cite{ChenFreedman2011, chambers2019minimum, Chambers2009, Borradaile2020, Grochow2018, DeyHouMandal2020}.


As our first contribution, we introduce a new type of barcode called  \emph{\hcb}. This barcode is obtained by tracking the birth and death of harmonic chains in an increasing filtration.  
Since there is always a \emph{unique} harmonic chain in a homology class, the \emph{harmonic representatives} in our {\hcb} immediately remove the second level of choices in a natural way. The main idea is to take the harmonic chain groups along an increasing filtration, and to observe that the groups grow or shrink due to the birth and death of harmonic chains, resulting in a zigzag module with inclusion maps. 
Like any zigzag module~\cite{CarlssonSilva2010}, the module of harmonic chain groups decomposes into interval modules, which then form the {\hcd}.
Moreover, since the maps in this zigzag module are essentially inclusions,
each representative for a bar in our {\hcb} consists of a single harmonic chain which is alive over the entire bar. We emphasize that our {\hcb} is distinct from the ordinary persistence barcode; see for example \Cref{fig:filtbarc,fig:twobc}. 

The ordinary persistence barcode is shown to be stable~\cite{Cohen-SteinerEdelsbrunnerHarer2005}, which is crucial for applications. The stability means that small changes in the data imply only small changes in the barcode. For our second contribution, we show that our {\hcb} is also stable in the same sense as the stability of persistence barcode~\cite{Cohen-SteinerEdelsbrunnerHarer2005}. 

Finally, we present an algorithm for computing 
the {\hcb} in $O(m^3)$ time for a filtration of size $m$,
matching the complexity of practical algorithms for ordinary persistence.

The interpretation of homological features is crucial for applications such as extracting hierarchical structures of amorphous solids~\cite{hiraoka2016hierarchical} and quantifying the growing branching architectures of plants~\cite{LiDuncanTopp2017}.
Such an interpretation is closely related to the existence of canonical chain representatives for a barcode. 
In our harmonic chain barcode, using orthogonality, each bar enjoys a canonical choice of representative (within a homology class) which \emph{lives exactly at the time-interval of the bar}. This arguably promises a more interpretable data feature. Therefore, we expect our {\hcbs} to enrich the family of topological descriptors in applications where ordinary persistence barcode is used, such as feature vectorization and machine learning. Note that just from the fact that our barcode is different from the ordinary persistence one cannot argue in favor or against our suggested barcode. Such a comparison is only meaningful in a specific domain of application and when performed using experimental evaluations. We refer to~\cite{GurnariGuzman-SaenzUtro2023} for an interesting pipeline of data analysis using extra properties of harmonic chains. \cite{GurnariGuzman-SaenzUtro2023} also contains a method of propagating the information from harmonic chains back to original data points which we find also relevant to our harmonic representatives.


%\Salman{this does not have a clear meaning and I think would trigger "anti-gemoetry" comments} Our work also provides strong evidence in a growing list of literature (e.g.,~\cite{BubenikHullPatel2020, BasuCox2022,ChacholskiGiuntiJin2023}) that rich geometric (not just topological) information can be recovered from a persistence filtration. 



% END AUTHOR FILE sec-introduction.tex
% BEGIN AUTHOR FILE sec-related-work.tex
\section{Related Work}
\label{sec:related-work}

Harmonic chains were first studied in the context of functions on graphs, where they were identified as the kernel of the Laplacian operator on graphs~\cite{Kirchhoff1847}. 
The graph Laplacian and its kernel are important tools in studying graph properties, see~\cite{Merris1994,MoharAlaviChartrand1991} for surveys. 
Eckmann~\cite{Eckmann1944} introduced the higher-order Laplacian for simplicial complexes, and proved the isomorphism of harmonic chains and homology. 
Guglielmi \etal~\cite{GuglielmiSavostianovTudisco2023} studied the stability of higher-order Laplacians. Horak and Jost~\cite{HorakJost2013} defined a weighted Laplacian for simplicial complexes. Already their theoretical results on Laplacian~\cite{HorakJost2013} anticipated the possibility of applications, as the harmonic chains are thought to contain important geometric information.
This has been validated by recent results that use curves of eigenvalues of Laplacians in a filtration in data analysis~\cite{ChenQiuWang2022, WangNguyenWei2020}. 
The Laplacian was applied to improve the mapper algorithm~\cite{MikePerea2019}, and for coarsening triangular meshes~\cite{KerosSubr2023}. 
The persistent Laplacian~\cite{MemoliWanWang2022} and its stability~\cite{LiuLiWu2023} is an active research area. 
Due to the close relation of harmonic chains and Laplacians, harmonic chains could find applications in areas that Laplacians have been used.
Computing reasonable representative cycles for persistent homology is also an active area of research. 
Here, usually an optimality criterion is imposed on cycles in a homology class to obtain a unique representative. For a single homology class, a number of works~\cite{Chen2010,DeyHiraniKrishnamoorthy2010, DeyHouMandal2020, Obayashi2018,Oleksiy2010, ChambersParsaSchreiber2022, ChenFreedman2011, DeyHiraniKrishnamoorthy2010, LeeChungChoi2019} consider different criteria for optimality of cycles. Hardness of computing optimal representatives has been studied by~\cite{ChenFreedman2011, chambers2019minimum, Chambers2009, ChambersParsaSchreiber2022, Borradaile2020, Grochow2018, DeyHouMandal2020}. For persistent homology, Dey et al.~\cite{DeyHouMandal2020} studied the hardness of choosing optimal cycles for persistence bars. 
Furthermore, De Gregorio et al.~\cite{DeGregorioGuerraScaramuccia2021} used harmonic cycles in a persistent homology setting to compute the persistence barcode. Lieutier~\cite{Lieutier} studied the harmonic chains in persistent homology classes, called persistent harmonic forms.

\para{Relation to the work of Basu and Cox.}~The most relevant work to ours is the inspiring work of Basu and Cox~\cite{BasuCox2022}. % , which has been recently used to study multi-omics data~\cite{GurnariGuzman-SaenzUtro2023}. 
Basu and Cox had a similar goal as ours, namely, to associate geometric information to each bar in order to obtain a more interpretable data feature. To that end, they introduced the notion of \emph{harmonic persistent barcode}, by associating a subspace of harmonic chains to each bar in the \textit{ordinary} persistence barcode. When the multiplicity of the bar is 1, the subspace is 1-dimensional. %This space is also defined for ordinary homology, denoted $M^{s,t}_p / N^{s,t}_p$, see~\cite[Section 3]{BasuCox2022}. %\Bei{What are $M$ and $N$?}
This is the space of harmonic chains that are born at $s$ and die entering $t$. In general, this is a quotient subspace of harmonic chains at time $s$. Using orthogonality %   from the given filtration of complexes, and analogously obtained a subquotient ${\cal M}^{s,t}_p / {\cal N}^{s,t}_p$. They used 
they can represent this subquotient as a subspace of the harmonic cycles at time $s$. As a result, they successfully assign a canonical harmonic cycle to a bar in the ordinary persistence diagram. We note that, in general, a bar in the persistence barcode cannot be represented using a single harmonic chain that \textit{remains harmonic during the lifetime of the bar}. The harmonic cycle associated to a bar using the Basu-Cox approach is the initial harmonic representative of the bar, that is, the harmonic cycle that represents the bar at its birth. Basu and Cox also used the \emph{terminal harmonic cycle} in some of their arguments, thus showing that the choice is not entirely canonical. However, Basu and Cox proved significant properties of the initial cycles in terms of what they call \emph{relative essential content}. The novelty of our result is that, in contrast to~\cite{BasuCox2022}, we define a new barcode distinct from the ordinary persistence barcode, in which each bar has a harmonic cycle associated with it which is harmonic during the lifetime of the bar. 
Basu and Cox also proved stability for their harmonic persistent barcode, by considering subspaces as points of a Grassmannian manifold and measuring distances in the Grassmannian. Such a distance quantifies the angles between subspaces, whereas our notions of stability are stronger in the sense that they use the classical bottleneck distance analogous to the ordinary persistence homology.
G\"{u}len et al.~\cite{GulenMemoliWan2024} studied a method that permits the construction of stable persistence diagrams that are equipped with a canonical choice of representative cycles for filtrations over arbitrary finite posets. If the underlying poset is a finite subset of the real line, using persistent Laplacians, they obtained harmonic cycles connected (through a certain isomorphism) to those identified by Basu and Cox. 








% END AUTHOR FILE sec-related-work.tex
% BEGIN AUTHOR FILE sec-fullversion-background.tex
\section{Background}
\label{sec:background}

In this section, we review the notion of harmonic chains and persistent homology. 
Homology and cohomology are defined with  real coefficients $\mathbb{R}$ (instead of $\mathbb{Z}_2$);
see~\cref{sec:homology} for a review of 
% standard notions of 
homology and cohomology.
% with real coefficients  $\mathbb{R}$  (instead of $\mathbb{Z}_2$).
Let $K$ be a simplicial complex and $p$ the homology dimension (or equivalently, homology degree). 
$C_p(K)$, $Z_p(K)$, and $H_p(K)$ denote the $p$-th chain group, cycle group, and homology group of $K$, whereas $C^p(K)$, $Z^p(K)$, and $H^p(K)$ denote the $p$-th cochain group, cocycle group, and cohomology group of $K$, respectively.
Groups across all dimensions are denoted as $C_*(K)$, $C^*(K)$, etc. 
We use $\partial$ and $\delta$ to denote boundary and coboundary operators, respectively. 

\subsection{Harmonic Cycles} 
Based on the standard notions of homology and cohomology  with $\Rspace$ coefficients (see~\cref{sec:homology}), we identify chains and cochains via duality, i.e.,~$C_p(K)= C^p(K)=\Rspace^{n_p}$, where $n_p$ is the number of $p$-simplices. 
Therefore, we can talk about coboundaries of cycles in $Z_p(K)$. 
We first introduce the notion of harmonic chains.  
\begin{definition}
\label{def:harmonic-cycle}
The \emph{$p$-th harmonic chain group} of $K$, denoted $\har_p(K)$, is the group of $p$-cycles that are also $p$-cocylces. 
Equivalently, $\har_p(K) := Z_p(K)\cap Z^p(K)$.
Each element in $\har_p(K)$ is called a \emph{harmonic $p$-chain}.
The harmonic chain group in all dimensions is the group $\har(K):=\bigoplus_p \har_p(K)$.
\end{definition}
% Following \cref{def:harmonic-cycle}, 
We sometimes use \emph{harmonic cycles} in place of \emph{harmonic chains} to emphasize the fact that harmonic chains are cycles.

\begin{lemma}[\cite{Eckmann1944}]
\label{lem:har-homo-iso}
$\har_p(K)$ is isomorphic to $H_p(K)$ and $H^p(K)$. Specifically, each homology and cohomology class has a unique harmonic cycle in it.
\end{lemma}

Harmonic cycles enjoy certain geometric properties. As an example, we mention the following \cref{prop:uniqueness}; see~\cite[Proposition 3]{DeSilvaMorozovVejdemoJohansson2011} for a proof. In other words, a harmonic cycle is the chain with the least squared-norm in a cohomology class.

\begin{proposition}[\cite{DeSilvaMorozovVejdemoJohansson2011}]
\label{prop:uniqueness}
Let $\alpha \in C^p(K)$ be a cochain. There is a unique solution $\bar{\alpha}$ to the least-squares minimization problem 
$\argmin_{\bar{\alpha}} \{ || \bar{\alpha}||^2 \mid \exists \gamma \in C^{p-1}(K);  \alpha = \bar{\alpha}+\delta \gamma \}.$ 
Moreover, $\bar{\alpha}$ is characterized by the relation $\partial {\bar{\alpha}}=0$.
\end{proposition}

There is an alternative definition of harmonic cycles. Consider the natural inner product on $C_p(K)$ given by $\langle\sigma_i , \sigma_j\rangle = \delta_{i,j}$. The harmonic chain group can be defined as $\har_p(K) = Z_p(K) \cap B_p(K)^\perp$. With this definition, the isomorphism of~\cref{lem:har-homo-iso} is realized by a  map that sends $z+B_p(K)$ to its projection to $B_p(K)^\perp$. In addition, the harmonic cycles satisfy 
$\har_p(K) = \ker(\partial_p) \cap \ker(\partial_{p+1}^\perp).$
For a short proof of the equivalence among the definitions of harmonic cycles above, see~\cite{BasuCox2022}. 
Importantly, our algorithm relies on the fact that harmonic cycles are cocycles whose boundaries are zero. 

\subsection{Ordinary Persistence}
\label{sec:persistence}

Ordinary persistent homology takes a \emph{filtration} of a simplicial complex $K$ as input. A \emph{continuous}  filtration $F$  assigns to each $r\in \Rspace$ a subcomplex $K_r \subseteq K$ such that  $K_r \subseteq K_s$ for $r \leq s$. 
Since $K$ is finite, there are finitely many  $t_0,t_1, \ldots, t_m \in \Rspace$ where $K_{t_i}$ changes. 
 We then abuse the notations slightly by
letting $K_i := K_{t_i}$ and have a \emph{discrete} form of $F$,
\begin{equation}
\label{eq:fcomplex}
    \filt:\emptyset = K_0 \hookrightarrow K_1 \hookrightarrow \cdots \hookrightarrow K_{m-1} \hookrightarrow K_m=K,
\end{equation}
where
each $K_i \hookrightarrow K_{i+1}$ is an inclusion. 
% Thus our filtrations are discretized using the real numbers $\{ t_i\}$ which we usually suppress in notations.
Unless stated otherwise, we assume that $\filt$ is \emph{simplex-wise}, i.e., each 
two 
% consecutive complexes 
$K_i$ and $K_{i+1}$ differ by at most a single simplex. 
% In this paper, 
For simplicity of the exposition,
complexes in $\filt$ sometimes are
subscripted by real-valued ``timestamps'' of the form $K_{t_i}$ (e.g., in~\cref{sec:stability}) or subscripted by integers of the form $K_i$ (e.g., in~\cref{sec:algorithm}), which  should not cause any confusions.
Applying homology functor to~\Cref{eq:fcomplex}, we obtain a sequence of homology groups and connecting linear maps (homomorphisms), forming a \emph{persistence module}:
\begin{equation}\label{eq:fhomology}
\Mcal: H_p(K_0) \to H_p(K_1) \to  \cdots \to H_p(K_{m-1}) \to H_p(K_m).   
\end{equation} 
For $s\leq t$, let $f^{s,t}_p: H_p(K_s) \rightarrow H_p(K_t)$ denote the map induced on the $p$-th homology by  inclusion.
% $K_s \subseteq K_t$. 
The image of the map, $f^{s,t}_p(H_p(K_s)) \subseteq H_p(K_t)$, is called the $p$-th \emph{$(s,t)$-persistent homology group}, denoted $H_p^{s,t}$. The group $H_p^{s,t}$
% , which in our case is also a vector space, 
consists of homology classes which exist in $K_s$ and survive until $K_t$. The dimensions of these vector spaces are the \emph{persistent Betti numbers}, denoted $\beta^{s,t}_p$.
An \emph{interval module}, denoted $I=I[b,d)$, is a persistence module of the form
\[ 0 \rightarrow \cdots \rightarrow 0 \rightarrow \Rspace \rightarrow \cdots \rightarrow \Rspace \rightarrow 0 \rightarrow \cdots \rightarrow 0.\]
In the above, $\Rspace$ is generated by a homology class and the connecting homomorphisms map generator to generator. 
We have $I_r = \Rspace$
for $b \leq r < d$ and $I_r=0$ for other $r$. 
Any persistence module can be decomposed into a collection of interval modules in a unique way \cite{LuoHenselman-Petrusek2023}. 
The collection of $[b,d)$ for all interval modules is called the \emph{persistence barcode}. 
When plotted as points in an extended plane, the result is the equivalent \emph{persistence diagram}. 
 
\para{Stability of Persistence Diagram/Barcode.}
The stability of persistence diagrams (or barcodes) is a crucial property for applications. It says that small changes in data lead to small changes in the persistence diagrams. 
We only review the stability for sublevel-set filtrations here; see \cref{subsec:ordinaryinterleaving} for more on stability of persistence modules.
Let $D$ and $D'$ denote two persistence diagrams. Recall that a persistence diagram is a multi-set of points in the extended plane (each of which is a birth-death pair) which also contains all points on the diagonal.
The \emph{bottleneck distance} of $D,D'$ is defined as 
\[
d_B(D, D') = \text{inf}_\gamma \text{sup}_{p \in D} || p - \gamma(p) ||_\infty,
\]
where $\gamma$ ranges over all bijections between $D$ and $D'$ and $||\cdot ||_\infty$ is the largest absolute value of differences of the points' coordinates.
  
A function $\tilde{f} : |K| \rightarrow \Rspace$ is called \emph{simplex-wise linear} if it is linear on each simplex.
Let $K_r = \{ \sigma \in K\mid\forall x \in |\sigma|, \tilde{f}(x) \leq r \}$ be the \textit{sublevel set complex} at value $r\in \Rspace.$ The complexes $K_r$ and the inclusions between them form a \textit{sublevel set filtration} $\tilde{F}$ (which is not necessarily simplex-wise). We denote the persistence diagram of $\tilde{F}$ as $\Dgm(\tilde{f})$.
We refer to~\cite{Cohen-SteinerEdelsbrunnerHarer2005, ChazalCohenSteinerGlisse2009} for proof of the following \Cref{theorem:persistence-stability}. See~\cite{Cohen-SteinerEdelsbrunnerHarer2005, ChazalCohenSteinerGlisse2009} for the stability of ordinary persistence.

\begin{theorem}
%\label{theorem:persistence-stability}
    Let $\tilde{f}, \tilde{g}: |K| \xrightarrow{} \mathbb{R}$ be simplex-wise linear functions. Then $$ d_B( \Dgm(\tilde{f}), \Dgm(\tilde{g})) \leq ||\tilde{f} - \tilde{g} ||_\infty.$$
\end{theorem}


 


%We also notice that in the above theorem we can replace $d_{CI}$ in place of $d_I$, since the existence of an $\eps$-chain-interleaving implies the existence of an $\eps$-interleaving on the homology level. 
% The resulting theorem would be weaker though.



\begin{comment}
 
Let $F$ and $G$ be two filtrations over the complexes $K$ and $K'$ respectively, and let the maps connecting the complexes in $F$ and $G$ be $f^{s,t}$ and $g^{s,t}$ respectively, for  $s\leq t \in \Rspace$. Let $C(F)$ and $C(G)$ be the corresponding filtrations of chain groups, and $H(F)$ and $H(G)$ the corresponding persistence modules. With an abuse of notation, we denote the maps induced on chain groups and homology groups also by $f^{s,t}$ and $g^{s,t}$ respectively, where $s\leq t \in \Rspace$.

\begin{definition}
\label{def:chain-interleaving}
Let $F$ and $G$ be two filtrations over  $K$ and $K'$ respectively. 
An \emph{$\eps$-chain-interleaving between $F$ and $G$} (or an $\eps$-interleaving at the chain level) is given by two sets of homomorphisms $\{\phi_\alpha : C(K_\alpha) \xrightarrow{} C(K'_{\alpha+\eps}) \}$ and $\{\psi_\alpha : C(K'_{\alpha}) \xrightarrow{} C(K_{\alpha+\eps}) \}$, such that 
\begin{enumerate}
    \item $\{ \phi_\alpha \}$ and $\{\psi_\alpha \}$ commute with the maps in the filtration, that is, for all $\alpha,t \in \Rspace$, $g^{\alpha+\eps,\alpha+\eps + t} \phi_\alpha = \phi_{\alpha+t} f^{\alpha, \alpha+t}$ and $f^{\alpha+\eps,\alpha+\eps+t} \psi_\alpha = \psi_{\alpha+t} g^{\alpha, \alpha+t}$; 
    \item The following diagrams commute: 
    \begin{equation}
      \begin{tikzcd}
  C(K_\alpha)
    \ar[r,"f^{\alpha, \alpha+\eps}"]
    \ar[dr, "\phi_\alpha", very near start, outer sep = -2pt]
  & C(K_{\alpha+\eps})
    \ar[r, "f^{\alpha+\eps, \alpha+2\eps}"]
    \ar[dr, "\phi_{\alpha+\eps}", very near start, outer sep = -2pt]
  & C(K_{\alpha+2\eps})
  \\
  C(K'_\alpha)
    \ar[r, swap,"g^{\alpha, \alpha+\eps}"]
    \ar[ur, crossing over, "\psi_\alpha"', very near start, outer sep = -2pt]
  & C(K'_{\alpha+\eps})
    \ar[r, swap,"g^{\alpha+\eps, \alpha+2\eps}"]
    \ar[ur, crossing over, "\psi_{\alpha+\eps}"', very near start, outer sep = -2pt]
  & C(K'_{\alpha+2\eps}).
  \end{tikzcd}
    \end{equation}
\end{enumerate}
\noindent The \emph{chain interleaving distance} is defined to be 
\begin{equation}
d_{CI}(F, G) := \inf \{\eps\geq 0 \mid \text{there exists an $\eps$-chain interleaving between } F\; \text{and}\; G\}.
\end{equation}
\end{definition}

We notice that the standard notion of $\eps$-interleaving~\cite{ChazalCohenSteinerGlisse2009}, denoted $d_I(F,G)$, is analogous to our definition above. However, it is defined on the persistence modules $H(F)$ and $H(G)$ on the homology level~\cite{ChazalCohenSteinerGlisse2009}, rather than the filtration of chain groups.  
Filtrations on the chain level arising from sublevel-sets of a function are ones that interest us in this paper.
% The filtration on the chain level of chains exists in the sublevel-set application that interests us.

The stability of persistence diagrams (or barcodes) is a crucial property for applications. It says that small changes in data lead to small changes in the persistence diagrams (PDs). While the interleaving distance provides a measure of change in data, we measure the distance between PDs using the bottleneck distance. Let $D$ and $D'$ denote two PDs. 
Recall that a PD is a multi-set of points in the extended plane (each of which is a birth-death pair) which also contains all points on the diagonal.
The \emph{bottleneck distance} is defined as 
\[
d_B(D, D') = \text{inf}_\gamma \text{sup}_{p \in D} || p - \gamma(p) ||_\infty,
\]
where $\gamma$ ranges over all bijections between $D$ and $D'$ and $||\cdot ||_\infty$ is the largest absolute value of differences of the points' coordinates.
% (or  the largest distortion of an endpoint 
% % when matched to the bar of $D'$ 
% when considering the PD points as bars).   
We refer to~\cite{ChazalCohenSteinerGlisse2009} for proof of the following \Cref{theorem:persistence-stability}.

\begin{theorem}
\label{theorem:persistence-stability}
    Let $F$ and $G$ be filtrations defined over (finite) complexes $K$ and $K'$, respectively. Then $$ d_B( Dgm(F), Dgm(G)) \leq d_I(F,G).$$
\end{theorem}

See~\cite{Cohen-SteinerEdelsbrunnerHarer2005, ChazalCohenSteinerGlisse2009} for more on the stability of ordinary persistence. We also notice that in the above theorem we can replace $d_{CI}$ in place of $d_I$, since the existence of an $\eps$-chain-interleaving implies the existence of an $\eps$-interleaving on the homology level. 
% The resulting theorem would be weaker though.

\end{comment}


\subsection{Zigzag Persistence}
\label{sec:zigzag}

We provide a brief overview of zigzag persistence; see~\cite{CarlssonSilva2010,CarlssonSilvaMorozov2009} for details.
A zigzag module 
\[\Mcal:V_0\lrarrowsp{g_0}V_1\lrarrowsp{g_1}\cdots\lrarrowsp{g_{k-1}}V_k\]
is a sequence of vector spaces connected by linear maps which could be forward or backward,
i.e., each $g_i$ could be $g_i:V_i\to V_{i+1}$ or $g_i:V_{i}\leftarrow V_{i+1}$.
The module $\Mcal$ decomposes into a direct sum of interval modules $I[b,d]$ of the form 
\[0 \longleftrightarrow \cdots \longleftrightarrow 0 \longleftrightarrow \Rspace \longleftrightarrow \cdots \longleftrightarrow \Rspace \longleftrightarrow 0 \cdots \longleftrightarrow 0 \]
with 1-dimensional vector spaces in the range $[b,d]$. 
The multi-set of intervals in the decomposition
defines the \emph{barcode} of $\Mcal$, denoted as $\barc(\Mcal)$.

\para{Conventions.}
% Since we work with chains and apply boundary and coboundary operators at different times in a filtration, we write $\delta_t(c)$ and $\partial_t(c)$ to mean the coboundary and boundary of the chain $c$ at time $t$, respectively, where we consider $c$ to be a chain in the final complex $K$ present at time $t$. 
In this paper, we may omit the subscript/dimension of a homology group if there is no danger of ambiguity. Moreover, we use the terms persistence barcode and persistence diagram interchangeably. 



% END AUTHOR FILE sec-fullversion-background.tex
% BEGIN AUTHOR FILE sec-barcode-zigzag.tex
\section{Harmonic Chain Barcodes and Representatives}
\label{sec:barcode-zigzag}

As reviewed in \cref{sec:persistence}, by directly taking the homology functor, an increasing filtration of simplicial complexes leads to an \emph{ordinary persistence module} consisting of homology groups~\cite{chazal2016structure,DeyWang2022}. These homology groups are connected by forward maps of the form $H(K_i) \xrightarrow{} H(K_{i+1})$. The ordinary persistence module then decomposes into interval  modules, which define the \emph{ordinary persistence barcode}. 
In this section, we show that the harmonic chain groups of all the complexes in a filtration constitute an abstract \emph{zigzag persistence module}  (see~\cref{sec:zigzag}). 
The main idea is that we take the harmonic chain groups along the filtration, where the groups could grow or shrink, resulting in a zigzag module. 
Moreover, the maps in this zigzag module are inclusions. 
Like any zigzag module~\cite{CarlssonSilva2010}, the module of harmonic chain groups decomposes into interval modules. These intervals form the {\hcd}, the main object of interest in this paper. 

Throughout the  section,
consider a simplex-wise filtration
\begin{equation*}
\filt: 
\emptyset
=K_0 \inctosp{\fsimp_0} K_1 \inctosp{\fsimp_1} 
\cdots \inctosp{\fsimp_{\filtcnt-1}} K_\filtcnt=K,
\end{equation*}
where each $K_{i+1}$ differs from $K_i$
by the addition of a simplex $\fsimp_i$. 
Recall that $K_i := K_{t_i}$ for $t_i\in\Real$ where $K_{t_i}$
is a complex from a \emph{continuous} filtration indexed over $\Real$; see \Cref{sec:persistence}. 

%\Tao{may need to mention the identification of cochains}

\begin{proposition}
\label{prop:arrow-dim-change}
For each inclusion $K_i \inctosp{\fsimp_i} K_{i+1}$ in $\filt$: 
\[\dim(\har(K_{i+1}))=\dim(\har(K_{i})) \pm 1.\]
\end{proposition}
\begin{proof}
This follows from \Cref{lem:har-homo-iso} and some well-known facts in persistence (see~\cite{DeyWang2022,EdelsbrunnerLetscherZomorodian2002}).
\end{proof}

\begin{definition}
\label{def:pos-neg-simp}
A simplex $\fsimp_i$ inserted in $\filt$ is called \emph{positive}
if $\dim(\har(K_{i+1}))=\dim(\har(K_{i}))+1,$ and \emph{negative}
if $\dim(\har(K_{i+1}))=\dim(\har(K_{i}))-1$.
\end{definition}

We describe how we connect the harmonic chain groups of complexes in $\filt$ by inclusions.
For each $K_i$ and each chain $c=\sum_{j=0}^{i-1}\aG_j\fsimp_j$ in $\Chn_p(K_i)$,
we identify $c$ as a chain $c=\sum_{j=0}^{\filtcnt-1}\aG_j\fsimp_j$ in $\Chn_p(K)$,
where $\aG_j=0$ for $j\geq i$.
This makes 
both 
$\Chn_p(K_i)$ and $\Chn_p(K_{i+1})$ a subspace of $\Chn_p(K)$. 
We then have the following inclusion:
\begin{equation}
\label{eqn:zyc-inc}
\Zyc_p(K_i) \incto \Zyc_p(K_{i+1}).
\end{equation}
Recall that $\har_p(K_i):=\Zyc_p(K_i)\cap \Zyc^p(K_i)$, which means that
$\har_p(K_i)\subseteq \Zyc_p(K_i) \subseteq \Chn_p(K)$.
Hence, we also identify each harmonic cycle in $\har_p(K_i)$ as a chain in $\Chn_p(K)$. 
We then observe in \cref{thm:harmonic-inclusion}
 a similar inclusion as in \Cref{eqn:zyc-inc} between any harmonic chain groups $\har_p(K_i)$ and $\har_p(K_{i+1})$,
with a possible flip on the direction.
\begin{theorem}
\label{thm:harmonic-inclusion}
For each arrow $K_i \inctosp{\fsimp_i} K_{i+1}$ in $\filt$, where $\fsimp_i$ is 
a $p$-simplex:  
\begin{itemize}
\item There is an inclusion $\har_p(K_i)\incto \har_p(K_{i+1})$
if $\fsimp_i$ is positive;
\item And there is an inclusion $\har_{p-1}(K_i)\bakincto\har_{p-1}(K_{i+1})$ if $\fsimp_i$ is negative.
\end{itemize}
In addition, in each case the harmonic chain groups in other dimensions remain unchanged, i.e., $\har_q(K_i) = \har_q(K_{i+1})$ for any other $q \notin \{p, p-1\}$.
\end{theorem}
\begin{proof}
The only harmonic chain groups that might change from $K_i$ to $K_{i+1}$ are 
those in dimension $p$ because cycle and cocycle groups in other dimensions  do not change.

First consider Case I that $\fsimp_i$ is positive.
Let $\iota: K_i \incto K_{i+1}$ be the inclusion map.  As noted above, we identify any $c \in \Chn_*(K_i)=\Chn^*(K_i)$ with $\iota_\sharp (c) \in \Chn_*(K_{i+1})=\Chn^*(K_{i+1})$.

\begin{description}
\item[Case I.1: $\har_p(K_i)\subseteq \har_p(K_{i+1})$.] 
Take $z \in \har_p(K_{i})$. 
Obviously, $\iota_\sharp(z)$ is a cycle in $K_{i+1}$.
Moreover, for any $c \in C_{p+1} (K_{i+1})$, $\delta (\iota_\sharp (z)) (c) = \iota_\sharp(z) (\partial c)$. Since $C_{p+1} (K_i)= C_{p+1} (K_{i+1})$, $c$ and hence $\partial c$ exist in $K_i$, meaning that $\iota_\sharp(z) (\partial c) = z (\partial c) = (\delta z) (c) = 0$. It follows that $\iota_\sharp(z)$ is a cocycle in $K_{i+1}$. Therefore, $\har_p(K_i) \subseteq \har_p (K_{i+1})$.

\item[Case I.2: $\har_{p-1}(K_i)= \har_{p-1}(K_{i+1})$.] 
Take $z \in \har_{p-1}(K_{i+1})$. 
Since $\fsimp_i$ is a positive $p$-simplex,
we have $Z_{p-1}(K_{i+1}) = Z_{p-1}(K_i)$
and $B_{p-1}(K_{i})=B_{p-1}(K_{i+1})$.
So $z \in Z_{p-1}(K_i)$
and we consider $z$ as a cochain in $K_i$.
Then, for any $c \in C_{p}(K_{i})$,  $\delta (z) (c) = z(\partial c) = 0$, with the last equality due to 
$\partial c\in B_{p-1}(K_{i})=B_{p-1}(K_{i+1})$.
Therefore, $\har_p(K_{i+1}) \subseteq \har_p(K_i)$.
Take $z \in \har_{p-1}(K_{i})$. For any $c \in C_{p}(K_{i+1})$,  $\delta (\iota_\sharp z) (c) = (\iota_\sharp z)(\partial c) = 0$, 
with the last equality due to 
$\partial c\in B_{p-1}(K_{i+1})=B_{p-1}(K_{i})$.
Therefore, $\har_{p-1}(K_i)= \har_{p-1}(K_{i+1})$.
\end{description}
    
Now consider Case II that $\fsimp_i$ is negative.

\begin{description}
\item[Case II.1: $\har_p(K_{i+1}) = \har_p(K_i)$.] 
Since $\fsimp_i$ is negative, $Z_p(K_{i+1}) = Z_p(K_i)$
and $B_p(K_{i+1})=B_p(K_{i})$.
The verification for this case is then the same as the verification
for Case I.2 with a shift on the homology degree.

% Take $z \in \har_p(K_{i+1})$. 
% Since $\fsimp_i$ is negative, $Z_p(K_{i+1}) = Z_p(K_i)$.
% So $z \in Z_p(K_i)$
% and we consider $z$ as a cochain in $K_i$.
% Then,
% for any $c \in C_{p+1}(K_{i})$,  $\delta (z) (c) = z(\partial c) = 0$, with the last equality due to 
% $\partial c\in B_p(K_{i})=B_p(K_{i+1})$.
% Therefore, $\har_p(K_{i+1}) \subseteq \har_p(K_i)$.
% Take $z \in \har_p(K_{i})$. For any $c \in C_{p+1}(K_{i+1})$,  $\delta (\iota_\sharp z) (c) = (\iota_\sharp z)(\partial c) = 0$, 
% with the last equality due to 
% $\partial c\in B_p(K_{i+1})=B_p(K_{i})$.  
% Therefore, $\har_p(K_{i+1}) = \har_p(K_i)$.
\item[Case II.2: $\har_{p-1}(K_{i+1}) \subseteq \har_{p-1}(K_i)$.]
Take $z \in \har_{p-1}(K_{i+1})$. We have $Z_{p-1}(K_{i+1}) = Z_{p-1}(K_i)$.
Therefore, $z \in Z_{p-1}(K_i)$
and we consider $z$ as a cochain in $K_i$. 
For any $c \in C_{p}(K_{i})$,  $\delta (z) (c) = z(\partial c) = 0$, 
with the last equality due to $\partial c\in B_{p-1}(K_i) \subseteq B_{p-1}(K_{i+1})$.
Therefore, $\har_{p-1}(K_{i+1}) \subseteq \har_{p-1}(K_i)$.\qedhere
\end{description}
\end{proof}

\begin{definition}[Harmonic chain barcode]
Consider the following \emph{harmonic zigzag module}
\begin{equation}
\label{eqn:har-zz-mod}
\har(\filt): \har(K_0) \leftrightarrow \har(K_1) \leftrightarrow
\cdots \leftrightarrow \har(K_\filtcnt),
\end{equation}
where the harmonic chain groups are connected by either forward inclusions {\rm(}e.g., $\har(K_i) \to \har(K_{i+1})${\rm)} or backward inclusions {\rm(}e.g., $\har(K_i) \leftarrow \har(K_{i+1})${\rm);}
see \Cref{thm:harmonic-inclusion}.
Define the \emph{\hcd} $\hbarc(\filt)$ of $\filt$ 
as the barcode of the zigzag module $\har(\filt)$,
that is, \[\hbarc(\filt):=\barc(\har(\filt)).\]
\end{definition}

In general, the {\hcb} is different from the ordinary persistence barcode for a filtration; see~\Cref{fig:twobc}.
In this paper, we also consider the zigzag module $\har_p(\filt)$ derived by taking $\har_p(K_i)$ on each $K_i$ in $\filt$.
While the $p$-th harmonic chain groups in $\har_p(\filt)$ are still connected by forward or backward inclusions, we may have $\har_p(K_i)=\har_p(K_{i+1})$ for two consecutive groups in $\har_p(\filt)$.
We therefore define the \emph{$p$-th {\hcb}} $\hbarc_p(\filt)$ of $\filt$ 
as the barcode of $\har_p(\filt)$,
i.e.,~$\hbarc_p(\filt):=\barc(\har_p(\filt))$.
Since $\har(\filt)=\bigoplus_p\har_p(\filt)$ (following from \Cref{thm:harmonic-inclusion}),
we have $\hbarc(\filt)=\bigsqcup_p\hbarc_p(\filt)$.


\para{Harmonic Representatives.}
In the rest of the section, we define harmonic representatives.
\begin{proposition}
Let $[b,d]$ be an interval in $\hbarc(\filt)$.
The inclusion $\har(K_{b-1}) \incto \har(K_b)$ in $\har(\filt)$
is forward.
Moreover, if $d<\filtcnt$,
then the inclusion $\har(K_{d}) \bakincto \har(K_{d+1})$ is backward.
\end{proposition}
\begin{proof}
This follows from \Cref{thm:harmonic-inclusion} and the fact that 
$\dim(\har(K_{b}))=\dim(\har(K_{b-1}))+1$ 
and
$\dim(\har(K_{d+1}))=\dim(\har(K_{d}))-1$.
% \Tao{probably need some modification}
\end{proof}

% \Tao{The rest is to show that the representative for each interval in this harmonic barcode consists
% of only one cycle.}
Representatives for general zigzag modules were introduced in~\cite{BendichPaul2013Haro,maria2015zigzag} (see also~\cite{zzrep}),
% The original definition~\cite{zzrep,maria2015zigzag} of a representative
% for a zigzag interval 
% % a {\red zigzag barcode}
which consist of a sequence of cycles
% satisfying several conditions 
for each interval.
However, since the harmonic chain groups are connected by inclusion maps in $\har(\filt)$,
we have:
\begin{proposition}\label{prop:single-rep}
Each representative for an interval in $\hbarc(\filt)$
contains a single harmonic cycle.
\end{proposition}

We then adapt the definition of representatives for
general zigzag modules~\cite{zzrep,maria2015zigzag} and 
define
harmonic representatives as follows:
\begin{definition}[Harmonic representative]\label{dfn:rep}
A \emph{harmonic $p$-representative}
{\rm(}or simply \emph{$p$-representative}{\rm)}
for an interval $[b,d]\in\hbarc_p(\filt)$ is a $p$-cycle 
$z\in\har_p(K_i)$ for $i\in[b,d]$ satisfying:
\begin{description}
\item[Birth condition:] $z$ is born in $\har_p(K_{b})$, i.e., $z\in\har_p(K_{b})\setminus\har_p(K_{b-1})$ {\rm(}notice that $b>0${\rm);}
\item[Death condition:] $z$ dies leaving $\har_p(K_{d})$, i.e., $z\in\har_p(K_{d})\setminus\har_p(K_{d+1})$ if $d<m$.
\end{description}
\end{definition}

Sometimes we relax the restriction of $[b,d]\in\hbarc_p(\filt)$ 
% in \Cref{dfn:rep}
and have a harmonic representative for an arbitrary integer interval $[b,d]\subseteq[0,\filtcnt]$.
\Cref{sec:rep-just} details the original definition of zigzag representatives
and justifies \Cref{prop:single-rep} and \Cref{dfn:rep}. 
% END AUTHOR FILE sec-barcode-zigzag.tex
% BEGIN AUTHOR FILE sec-barcode-zigzag-algo.tex
\section{A Cubic-Time Algorithm for Computing Harmonic Chain Barcodes}\label{sec:algorithm}
In this section, we propose an $O(\filtcnt^3)$ algorithm for computing the {\hcb}  and its harmonic representatives given a filtration containing $\filtcnt$ insertions.
We first overview the algorithm and then describe the implementation.
While there have been algorithms~\cite{CarlssonSilvaMorozov2009,zzrep,maria2015zigzag} for computing zigzag barcodes in the general case, these algorithms target zigzag modules induced from directly \emph{taking the homology functor on zigzag filtrations}.
In contrast, our algorithm targets the special type of zigzag module $\har(\filt)$, where harmonic chain groups are derived from an \emph{ordinary non-zigzag} filtration and
are connected by inclusions.

For simplicity, when describing the algorithm, complexes in a filtration are always 
indexed by integers instead of the real-valued timestamps (see~\Cref{sec:persistence}).
Again, we assume that the input is a simplex-wise filtration 
$\filt: 
\emptyset
=K_0 \inctosp{\fsimp_0} K_1 \inctosp{\fsimp_1} 
\cdots \inctosp{\fsimp_{\filtcnt-1}} K_\filtcnt=K,
$
where each $K_{i+1}$ differs from $K_i$ by the addition of a simplex $\fsimp_i$. 

\Cref{alg:frame} provides an overview: our algorithm computes the harmonic chain barcode by iteratively maintaining the harmonic representatives. 
It processes the insertions in $\filt$ one by one and finds \emph{pairings} of the \emph{birth indices} (starting points of the intervals) and \emph{death indices} (ending points of the intervals) to form intervals in $\hbarc(\filt)$. 
When we encounter a new birth index, it is initially \emph{unpaired};
when we encounter a new death index, an unpaired birth index is chosen to pair with the death index. 

The following definition helps present the algorithm.
\begin{definition}
For an interval $[b,i] \subseteq [0,\filtcnt]$,
a \emph{partial $p$-representative} for $[b,i]$ is a $p$-representative as in \Cref{dfn:rep}
by ignoring the death condition.
\end{definition}

\makeatletter
  \def\vhrulefill#1{\leavevmode\leaders\hrule\@height#1\hfill \kern\z@}
\makeatother

\newenvironment{algrl}
    {\vspace{5pt}
\par\nobreak\nobreak\noindent\vhrulefill{0.8pt}
\begin{algr}
\vspace{-5pt}
    }
    { 
\vspace{-3pt}    
\par\nobreak\nobreak\noindent\vhrulefill{0.8pt}
\end{algr}
\vspace{3pt}
    }

\begin{algrl}\label{alg:frame}
Let $U^p$ be the set of unpaired birth indices for each homology degree $p$, where $U^p=\emptyset$ initially. Before each iteration that processes the insertion $K_{i}\inctosp{\fsimp_i}K_{i+1}$, we maintain a $p$-cycle $z$ for each $b\in U^p$ that is a partial representative for the interval $[b,i]$.
We use partial representatives to determine a finalized representative when a birth index is paired with a death index (by making sure the death condition is satisfied).

When processing the insertion of a $p$-simplex $\fsimp_i$ via $K_{i}\inctosp{\fsimp_i}K_{i+1}$, we proceed as follows:
\begin{description}
    \item[If $\sG_i$ is positive:] 
    \begin{itemize}
        \item[]
        \item Since $\dim(\har_p(K_{i+1}))=\dim(\har_p(K_{i}))+1$ (by \Cref{thm:harmonic-inclusion}), we add a new birth index $i+1$ to $U^{p}$.
        \item Find a harmonic $p$-cycle $z\in\har_p(K_{i+1})\setminus\har_p(K_{i})$, and let $z$ be the partial representative for $i+1\in U^p$.
    \end{itemize}
    
    \item[If $\sG_i$ is negative:] 
        \begin{itemize}
        \item[]
        \item Since $\dim(\har_{p-1}(K_{i+1}))=\dim(\har_{p-1}(K_{i}))-1$, we have a new death index $i$.
        \item Let $U^{p-1}=\{b_j\mid j=1,\ldots,k\}$, where each $b_j$ has a partial $(p-1)$-representative $z_j$.
        \item Let $U'=\{b_j\in U^{p-1}\mid z_j\not\in\har_{p-1}(K_{i+1})\}$, i.e., $U'$ contains all  birth indices in $U^{p-1}$ whose partial representatives do not persist to $K_{i+1}$.
        \item Pair the \emph{smallest} (i.e., the ``oldest'') index $b_{j^*}\in U'$ with $i$ which forms a new interval $[b_{j^*},i]\in\hbarc_{p-1}(\filt)$. Assign $z_{j^*}$ as the representative for $[b_{j^*},i]\in\hbarc_{p-1}(\filt)$ and remove $i$ from $U^{p-1}$.
        \item Consider each $b_{j}\in U'\setminus\{b_{j^*}\}$. We have that $z_j\not\in\har_{p-1}(K_{i+1})$ because $\dG(z_j)$ becomes non-zero in $K_{i+1}$ ($\partial z_j$ is always zero after being born). Also, the only reason that  $\dG(z_j)\neq 0$ in $K_{i+1}$ is because  $\aG_j:=\dG(z_j)(\fsimp_i)=z_j(\partial\fsimp_i)\neq 0$. Let $\aG_*=z_{j^*}(\partial\fsimp_i)$. Update the partial representative for $b_j$ as $z_j:= z_j-(\aG_j/\aG_*)\cdot z_{j^*}$ so that $z_j$ now persists to $\har_{p-1}(K_{i+1})$.
\end{itemize}
\end{description}

After processing all the insertions, for each $b$ in each $U^p$ with a partial representative $z$, let $[b,\filtcnt]$ form an interval in $\har_p(\filt)$ with a representative $z$.
\end{algrl}

\begin{example*}
\Cref{fig:filtbarc} provides an example of how \Cref{alg:frame} computes {\hcb} $\har_1(\filt)$ and representatives for a filtration $\filt$, with the resulting barcode  in \Cref{fig:twobc}. For simplicity, $K_0, K_1, K_2$ are omitted and vertex insertions are ignored. 
In $K_4$, $K_5$, and $K_7$, new 1-cycles are born which are also harmonic cycles due to a lack of triangles. The partial representatives $z_4$, $z_5$, $z_7$ for the indices $U^1=\{4,5,7\}$ all persist till $K_7$. 
When the triangle $abe$ is inserted in $K_8$ (producing a death index $7$), $\dG(z_4)(abe)=1$, $\dG(z_5)(abe)=1$, and $\dG(z_7)(abe)=3$.
We pair $4\in U^1$ with $7$ and have $[4,7]\in\har_1(\filt)$ with a representative $z_4$.
We also let $z_5:=z_5-z_4$ and $z_7:=z_7-3z_4$, which now both persist to $\har_1(K_8)$.
The remaining steps are similar. Notice that when $abe$ is inserted, the ``alive'' bar represented by the boundary of $abe$ (i.e.,~$ab+be-ae$) is not killed.
This is a significant difference from ordinary persistence (which will kill the alive bar represented by $\partial(abe)$ when $abe$ is inserted).
\end{example*}

\begin{figure}[!ht]
\centering
\includegraphics[width=\linewidth]{figs/filtbarc.pdf}
\caption{An example of the computation of {\hcb} $\har_1(\filt)$ and representatives.}
\label{fig:filtbarc}
\end{figure}

\begin{figure}[!ht]
  \centering
  \includegraphics[width=0.75\linewidth]{twobc}
  \caption{Harmonic chain barcode and ordinary persistence barcode for the filtration in \Cref{fig:filtbarc}. Deviating from conventions in ordinary persistence, bars are drawn as closed integer intervals, e.g., $[7,7]$ in the ordinary barcode is killed by the addition of $abe$ in $K_8$.}
  \label{fig:twobc}
\end{figure}

\begin{remark}
\Cref{alg:frame} chooses the ``oldest'' birth, among the options, to pair with a death while the ordinary persistence chooses the ``youngest'' one~\cite{EdelsbrunnerLetscherZomorodian2002}.
A brief explanation is that a summation of representatives \emph{persists} to the next harmonic chain group in $\har(\filt)$, whereas a summation of representatives
in the ordinary persistence \emph{becomes trivial} in the next homology group (and cannot persist).
We refer to~\cite{zzrep,maria2015zigzag} for a formal definition of the different types of birth/death ends and how they impact the computation of persistence.
\end{remark}

The following (rephrased) proposition from~\cite{dey2021computing} helps prove the correctness of \Cref{alg:frame}:

\begin{proposition}[Proposition 9, \cite{dey2021computing}]
\label{prop:pn-paring-w-rep}
If a pairing $\pi$ of the birth and death indices in $\har(\filt)$ satisfies that each interval from $\pi$ has a harmonic representative, then $\pi$ induces the {\hcb} $\hbarc(\filt)$.
\end{proposition}

\begin{theorem}
\Cref{alg:frame} correctly computes the {\hcb} for filtration $\filt$.
\end{theorem}
\begin{proof}
We first show by induction that before \Cref{alg:frame} processes each $\fsimp_i$, the $p$-cycle $z$ maintained for any $b\in U^p$ is a partial $p$-representative for $[b,i]$.
Suppose that the claim is true before processing $\fsimp_i$.
If $\fsimp_i$ is positive, index $i+1$ is added to $U^p$ and we assign $z \in\har_p(K_{i+1})\setminus\har_p(K_i)$ to $i+1$, which is clearly a partial representative for $[i+1,i+1]$.
For an index $b$ previously in $U^p$ having a partial representative $z$, since $\har_p(K_i)\subseteq \har_p(K_{i+1})$ (\Cref{thm:harmonic-inclusion}), we have that $z$ is still a partial representative for the interval $[b,i+1]$.

In the case that $\fsimp_i$ is negative, whenever we update a representative $z_j$ as $z_j-(\aG_j/\aG_*)\cdot z_{j^*}$, we have that $b_{j^*}<b_j$ and hence the updated $z_j$ remains born in $\har_{p-1}(K_{b_j})$. Therefore, the new $z_j$ is still a partial representative for $[b,i]$, which is also a partial representative for $[b,i+1]$ because now $z_j\in\har_{p-1}(K_{i+1})$.

By \Cref{prop:pn-paring-w-rep}, we only need to show that each interval output by the algorithm admits a harmonic representative.
This follows from the algorithm. 
For example, whenever we output an interval $[b_{j^*},i]\in\hbarc_{p-1}(\filt)$ when $\fsimp_i$ is negative, we have $z_{j^*}\not\in\hbarc_{p-1}(K_{i+1})$, which ensures the death condition.
\end{proof}


% \input{sec-barcode-zigzag-algo-impl-short}
% BEGIN AUTHOR FILE sec-barcode-zigzag-algo-impl.tex
% \let\mbeg\(
% \let\mend\)
\let\mbeg\[
\let\mend\]


%%%%%% BELOW ARE START OF PARAGRAPH FOR SOCG PUT IN APPENDIX !!
% \section{Full Implementation Details of \Cref{alg:frame}}\label{sec:impl}
% We provide the full 
% implementation details of \Cref{alg:frame}
% in this section.
%  Some of the statements 
% already presented in \Cref{sec:algorithm}
% (e.g., \Cref{algr:mainbody-reduce}, \Cref{prop:mainbody-reduce}, and \Cref{thm:mainbody-m3})
% are repeated in this section to prevent jumping back and forth.
%%%%%% ABOVE ARE START OF PARAGRAPH FOR SOCG PUT IN APPENDIX !!

\subsection{Implementation}\label{sec:impl}
Before detailing the implementation,
we first have the following fact:
% (proof in \Cref{sec:pf-prop-dim-chng}):

\begin{proposition}\label{prop:dim-chng}
When  a $p$-simplex $\fsimp_i$ is added  from $K_i$ to $K_{i+1}$,
one has:
   \begin{itemize}
   \item If $\fsimp_i$ is positive, then 
   \begin{align*}
       &\dim(\Zyc^p(K_{i+1}))=\dim(\Zyc^p(K_{i}))+1\text{ and}\\
   &\dim(\Bnd^p(K_{i+1}))=\dim(\Bnd^p(K_{i})).
   \end{align*}

   \item If $\fsimp_i$ is negative, then
   \begin{align*}
   &\dim(\Zyc^p(K_{i+1}))=\dim(\Zyc^p(K_{i}))+1\text{, }\\
   &\dim(\Bnd^{p}(K_{i+1}))=\dim(\Bnd^{p}(K_{i}))+1\text{, and}\\
   &\dim(\Zyc^{p-1}(K_{i+1}))=\dim(\Zyc^{p-1}(K_{i}))-1.
   \end{align*}
   \end{itemize}
Furthermore, in each case, all other cocycle or coboundary  groups not mentioned above stay the same
from $K_i$ to $K_{i+1}$.
\end{proposition}
\begin{proof}
We first show that $\dim(\Zyc^p(K_{i+1}))=\dim(\Zyc^p(K_{i}))+1$
whether $\fsimp_i$ is positive or negative.
% Let $\iG_\sharp:\Chn_p(K_{i})\to\Chn_p(K_{i+1})$
% be the inclusion of the two chain groups
Let \mbeg\iG^\sharp:\Chn^p(K_{i+1})\to\Chn^p(K_{i})\mend
be the cochain map induced by inclusion.
% the {\red dual} of $\iG_\sharp$.
Clearly, we have  $\iG^\sharp(\Zyc^p(K_{i+1}))\subseteq\Zyc^p(K_{i})$~\cite{hatcher2002algebraic}.
Let \mbeg\jmath^\sharp:\Zyc^p(K_{i+1})\to\Zyc^p(K_{i})\mend be the restriction of $\iG^\sharp$.
% to the domain $\Zyc^p(K_{i+1})$.
We show that $\jmath^\sharp$ is surjective
% (somehow obvious)
and the kernel of $\jmath^\sharp$ is 1-dimensional.
First, for any cocycle $\phi\in\Zyc^p(K_{i})$,
we define a cocycle $\phi'\in\Zyc^p(K_{i+1})$
where $\phi'(\fsimp_i)=0$
and values of $\phi'$ on other $p$-simplices are
the same as $\phi$.
We have that $\jmath^\sharp(\phi')=\phi$ and so $\jmath^\sharp$
is surjective.
Moreover, let $\hat{\fsimp}_i\in\Zyc^p(K_{i+1})$ be 
the dual of $\fsimp_i$.
% a cocycle
% such that $\hat{\sG}_i(\fsimp_i)=1$ and values of $\hat{\sG}_i$
% on other $p$-simplices are zero.
% \Tao{the dual should defined before?}
We have that $\hat{\fsimp}_i$ generates 
the kernel of $\jmath^\sharp$.
% \Tao{more explanation on this?}

We also notice that the coboundary maps \mbeg\dG^{p-2}: \Chn^{p-2}(K_i)\to \Chn^{p-1}(K_i)\mend
   and \mbeg\dG^{p-2}: \Chn^{p-2}(K_{i+1})\to \Chn^{p-1}(K_{i+1})\mend
   are the same, and so $\Bnd^{p-1}(K_{i+1})=\Bnd^{p-1}(K_{i})$.

Suppose that $\fsimp_i$ is positive.
We have  $\dim(\Hm^p(K_{i+1}))=\dim(\Hm^p(K_{i}))+1$.
% (see Lemma~\ref{lem:har-homo-iso}).
Since 
\begin{align*}
&\dim(\Hm^p(K_{i}))=\dim(\Zyc^p(K_{i}))-\dim(\Bnd^p(K_{i})),\\
&\dim(\Hm^p(K_{i+1}))=\dim(\Zyc^p(K_{i+1}))-\dim(\Bnd^p(K_{i+1})),
\text{ and }\\
&\dim(\Zyc^p(K_{i+1}))=\dim(\Zyc^p(K_{i}))+1,
\end{align*}
we have $\dim(\Bnd^p(K_{i+1}))=\dim(\Bnd^p(K_{i}))$.
Similarly, since
\begin{align*}
&\Hm^{p-1}(K_{i+1})\approx\Hm^{p-1}(K_{i}),\\
&\Bnd^{p-1}(K_{i+1})=\Bnd^{p-1}(K_{i}),
\text{ and }\\
&\Chn^{p-1}(K_{i+1})=\Chn^{p-1}(K_{i}),
\end{align*}
we have $\Zyc^{p-1}(K_{i+1})=\Zyc^{p-1}(K_{i})$.

% Moreover,
When $\fsimp_i$ is negative,
the statements regarding the dimensions of the 
cocycle and coboundary groups 
% as stated 
in the proposition
can be similarly obtained.

To see that other cocycle or coboundary groups do not change, 
we provide arguments for the groups in degree $p+1$ 
and the justification for groups in other degrees is omitted.
We first have $\Zyc^{p+1}(K_{i+1})=\Zyc^{p+1}(K_{i})$
because the two 
maps $\dG^{p+1}: \Chn^{p+1}(K_i)\to \Chn^{p+2}(K_i)$ and
$\dG^{p+1}: \Chn^{p+1}(K_{i+1})\to \Chn^{p+2}(K_{i+1})$ are the same.
Moreover, 
% since $\Hm^{p+1}(K_{i+1})\approx\Hm^{p+1}(K_{i})$
% and $\Chn^{p+1}(K_{i+1})=\Chn^{p+1}(K_{i})$,
we have $\Bnd^{p+1}(K_{i+1})=\Bnd^{p+1}(K_{i})$
because $\sG_i$ has no cofaces in $K_{i+1}$.
\end{proof}


\newcommand{\ldim}{{p}}
\para{Pseudo-cochains.}
We introduce the notion of {pseudo-cochains}
which are more ``concise'' versions
of  cochains.
% which completely determine the coboundaries.
% of a complex.
% As  seen later,
These
pseudo-cochains help maintain a coboundary basis
in our implementation.
Proofs of 
% the relevant propositions 
\Cref{prop:cobound-by-bound,prop:pseudocob-is-cob,prop:pseudo-coc-nat}
% concerning the pseudo-cochains
follow directly from
definitions
and are omitted.
% \Cref{prop:cobound-by-bound} 
% We first present 
\Cref{prop:cobound-by-bound}
% the following proposition
% {\red(proof follows directly from definitions)}  
says that 
the coboundary of a $\ldim$-cochain $\phi$
is completely determined by how $\phi$ maps the $\ldim$-boundaries 
to  real values:
\begin{proposition}\label{prop:cobound-by-bound}
Let
$\phi:\Chn_\ldim(K_i)\to\Real$ 
and $\psi:\Chn_\ldim(K_i)\to\Real$ be two $\ldim$-cochains  of $K_i$.
Also,
let $\bar{\phi}:\Bnd_\ldim(K_i)\to\Real$ and $\bar{\psi}:\Bnd_\ldim(K_i)\to\Real$ be restriction of $\phi$ 
and $\psi$
to $\Bnd_\ldim(K_i)$ respectively.
One has that  $\dG(\phi)=\dG(\psi)$
if and only if $\bar{\phi}=\bar{\psi}$.
\end{proposition}


\begin{definition}
Define a \emph{$\ldim$-pseudo-cochain} $\psi$ in $K_i$ as 
a map $\psi:\Bnd_\ldim(K_i)\to\Real$.
Define the \emph{coboundary} of $\psi$
as the $(\ldim+1)$-cochain $\psi\circ\bar{\partial}$,
where $\bar{\partial}:\Chn_{\ldim+1}(K_i)\to\Bnd_\ldim(K_i)$
is the restriction of $\partial$ to its image.
We abuse the notation slightly and denote 
the coboundary of $\psi$ as $\dG(\psi)$.
\end{definition}

\Cref{prop:pseudocob-is-cob} 
% {\red(proof follows directly from definitions)} 
says that coboundaries of
a complex  are completely determined by the pseudo-cochains:

\begin{proposition}\label{prop:pseudocob-is-cob}
The space of coboundaries of all 
$\ldim$-pseudo-cochains of $K_i$ equals $\Bnd^{p+1}(K_i)$.
\end{proposition}
% \begin{proposition}\label{prop:pseudocob-is-cob}
% Let  $\psi$ be a {$\ldim$-pseudo-cochain} of $K_i$.
% Then there exists a $\ldim$-cochain $\phi$ of $K_i$
% s.t.\ $\dG(\phi)=\dG(\psi)$.
% Also, for an arbitrary 
% $\ldim$-cochain $\phi'$ of $K_i$,
% let $\bar{\phi'}:\Bnd_\ldim({K_i})\to\Real$
% be the restriction of $\phi'$ to $\Bnd_\ldim({K_i})$.
% Then $\dG(\phi')=\dG(\bar{\phi'})$.
% \end{proposition}

\begin{definition}
Given
a basis $\Rcal=\{\zG_1,\ldots,\zG_r\}$ of $\Bnd_\ldim(K_i)$
and
a column vector $\vec{v}=(\aG_1,\ldots,\aG_r)^T$,
we say that  $\vec{v}$ \emph{encodes} 
the $\ldim$-pseudo-cochain $\psi:\Bnd_\ldim(K_i)\to\Real$
w.r.t\ 
% (with respect to)
the basis $\Rcal$
if $\psi(\zG_j)=\aG_j$ for each $j$.
\end{definition}

\Cref{prop:pseudo-coc-nat} 
says that summations and scalar multiplications
on the column vectors encoding 
$\ldim$-pseudo-cochains
% are natural with respect to 
commute with
the coboundary operator:

\begin{proposition}\label{prop:pseudo-coc-nat}
Fix
a basis $\Rcal=\{\zG_1,\ldots,\zG_r\}$ of $\Bnd_\ldim(K_i)$.
Let
$\vec{u}$, $\vec{v}$
be two column vectors encoding 
$\ldim$-pseudo-cochains $\psi$, $\psi'$
w.r.t $\Rcal$ respectively.
For an $\aG\in\Real$,
one has that $\vec{u}+\aG\vec{v}$ encodes
the $\ldim$-pseudo-cochains $\psi+\aG\cdot\psi'$
w.r.t $\Rcal$
and $\dG(\psi+\aG\cdot\psi')=\dG(\psi)+\aG\cdot\dG(\psi')$.
\end{proposition}

In this paper, when a basis
$\Rcal$ of $\Bnd_\ldim(K_i)$ is clear,
we abuse the notations and let $\dG(\vec{v}):=\dG(\psi)$,
where
$\vec{v}$ is a vector
encoding
a $\ldim$-pseudo-cochain $\psi$ w.r.t $\Rcal$.

\para{Representation of chains and cochains.}
% First, 
% In our implementation,
We represent
chains   as matrix columns
and make no distinction between a  column vector
and the chain it represents.
Specifically, 
% taking a fixed  orientation
% for each simplex $\fsimp_i$ in $\filt$,
with $\Chn_*(K_i)\subseteq\Chn_*(K)$
and a natural basis $\{\fsimp_0,\fsimp_1,\ldots,\fsimp_{\filtcnt-1}\}$
for $\Chn_*(K)$,
each chain \mbeg c=\sum_{k=0}^{i-1}\aG_k\fsimp_k\in\Chn_*(K_i)\mend
is represented as a column vector  \mbeg(\aG_0,\aG_1,\ldots,\aG_{\filtcnt-1})^T,\mend
where $\aG_k=0$ for each $k\geq i$.
% We also make no distinction between a matrix column
% and the chain it represents.
For a matrix $\mathtt{M}$, we denote its $j$-th
column as $\mathtt{M}[j]$.
For a chain $z$ as a column vector,
we  denote its $k$-th entry as $z[k]$.
Hence, the entry in the $k$-th row and $j$-th column 
in $\mathtt{M}$ is denoted as $\mathtt{M}[j][k]$.
Moreover,
we let the \emph{pivot} of $\mathtt{M}[j]$
be the index of the lowest non-zero entry in 
$\mathtt{M}[j]$
which is denoted as $\pivot(\mathtt{M}[j])$.
% In our implementation,
% Notice that

Cochains are similarly represented as matrix columns
by taking the dual basis \mbeg\{\hat{\fsimp}_0,\hat{\fsimp}_1,\ldots,\hat{\fsimp}_{\filtcnt-1}\}\mend
for $\Chn^*(K)$.
Notice that since chains and cochains
are  identified, 
a matrix column \mbeg(\aG_0,\aG_1,\ldots,\aG_{\filtcnt-1})^T\mend
can represent \emph{both} the chain $\sum_{k=0}^{\filtcnt-1}\aG_k\fsimp_k$ and the cochain
$\sum_{k=0}^{\filtcnt-1}\aG_k\hat{\fsimp}_k$.
Moreover, due to the identification,
 a formal sum \mbeg\sum_{k=0}^{\filtcnt-1}\aG_k\hat{\fsimp}_k\mend
% (resp.\ )
can also denote a chain,
% This is for the ease of the presentation and 
which should not cause any confusions.


\para{Matrices maintained.}
% Our algorithm processes each insertion 
% $K_{i}\inctosp{\fsimp_i}K_{i+1}$
% in $\filt$ one by one 
% % starting from 
% starting from $i=0$.
% As mentioned, our algorithm processes each insertion of $\fsimp_i$ in $\filt$ starting from $\fsimp_0$.
The only step in \Cref{alg:frame}
that cannot be easily implemented
is finding a new harmonic cycle born in $\har_p(K_{i+1})$
when 
% the $p$-simplex 
$\fsimp_i$ 
is positive.
We can do this naively by computing a basis for the  kernel of $\delta$ in $K_{i+1}$, and check all   cocycles in the basis to see if they are independent of the existing set of harmonic representatives
maintained for $K_i$. This, however, requires at least $\Theta(m^3)$ per simplex insertion.
% {\red In this case, the dimension of the harmonic chain group increases and we need to find a single harmonic cycle independent of the previous harmonic group. We can do that naively by computing a basis for the nullity of $\delta$ at the new complex, and check all basis elements to see if they are independent of our set of harmonic representatives. This, however, is at least $\Theta(m^3)$ per simplex insertion.}


For a more efficient implementation,
before  iteration $i$ which processes the insertion 
$K_{i}\inctosp{\fsimp_i} K_{i+1}$,
we maintain the following matrices
for each homology degree $p$:
% It turns out that a more efficient implementation
% entails maintaining several matrices
% serving different purposes
% for each homology degree $p$.
% Specifically,
% before an iteration $i$ which processes the insertion 
% $K_{i}\incto K_{i+1}$,
% we maintain the following matrices:
\begin{enumerate}

    \item $\Harmat^{p}$: Columns in  $\Harmat^{p}$
    are
    partial $p$-representatives 
    for all indices in $U^p$,
    % for all unpaired birth indices
% in {\red degree} $p$
% in Algorithm~\ref{alg:frame}.
% Notice that 
 % these columns
 which
 form a \emph{harmonic cycle basis} for $\har_p(K_i)$;
 see \Cref{sec:harmat-basis} for justification.
 % in each iteration 
 % of \Cref{alg:frame}
% as we scan through the filtration $\filt$
% ({\red see Proposition~XX}).

\item $\Rmat^p$: Columns in $\Rmat^p$ form a \emph{boundary basis} with distinct pivots for $\Bnd_p(K_i)$.
% in each iteration.
Notice that $\Rmat^p$ is indeed the ``reduced'' matrix 
in the
ordinary persistence algorithm~\cite{cohen2006vines}
with only $p$-chains
and without zero columns.

\item $\Cobmat^p$: Columns in $\Cobmat^p$ form a \emph{coboundary basis}
for $\Bnd^p(K_i)$.
% in each iteration.

% Notice that columns in $\Harmat^{p}$ and 
% $\Cobmat^p$
% altogether form a \emph{cocycle basis} for  $\Zyc^p(K_i)$.

\item 
$\Chnmat^p$: Columns in $\Chnmat^p$
encode $p$-pseudo-cochains in $K_i$ w.r.t\ the basis $\Rmat^p$
% $\{\}$
% such that
whose coboundaries are columns in
$\Cobmat^{p+1}$.
% , i.e.,
% the two matrices  
More specifically,
$\Chnmat^p$, $\Cobmat^{p+1}$ have the same number of columns
and for each $j$,
$\Cobmat^{p+1}[j]=\dG(\psi_j)$ 
where $\psi_j$ is the $p$-pseudo-cochain
encoded by $\Chnmat^p[j]$
w.r.t\ 
% the basis 
$\Rmat^p$.

\item $\Bndmat^p$: Columns in $\Bndmat^p$ are boundaries of the
coboundaries in $\Cobmat^{p+1}$,
% of the columns in $\Chnmat^{p}$,
i.e., 
% the two matrices 
 $\Bndmat^p$, $\Cobmat^{p+1}$
have the same number of columns and
$\Bndmat^p[j]=\partial(\Cobmat^{p+1}[j])$ for each $j$.
We further ensure that columns in $\Bndmat^p$ have distinct pivots
so that the reduction process in \Cref{algr:reduce} 
% (presented later)
% ; {\red essential for the finding new-born harmonic cycles}) 
can be done in $O(\filtcnt^2)$ time.
\end{enumerate}

% \begin{remark}
In summary,
columns in   $\Chnmat^{p-1}$, $\Cobmat^{p}$,  and  $\Bndmat^{p-1}$
have one-to-one correspondence
as follows:
\[
\Chnmat^{p-1}[j]\xmapsto{\ \dG\ }\Cobmat^{p}[j]\xmapsto{\ \partial\ }\Bndmat^{p-1}[j].\]
% \Tao{say something more about why do we need all these matrices?}
% \end{remark}

% \begin{remark}
Other than $\Harmat^p$ which is already in \Cref{alg:frame},
the rationale for maintaining the matrices is as follows
(assuming $\fsimp_i$ is a $p$-simplex):
\begin{itemize}
    \item We use $\Bndmat^{p-1}$ to find a linear combination $\phi$
    of  $\hat{\fsimp}_i$ and cocycles in $\Cobmat^{p}$ 
    such that $\partial\phi=0$
    (see \Cref{algr:reduce}).
    % whose boundary is zero.
    The cocycle $\phi$ 
    % with $\partial z=0$ 
    is then a new harmonic $p$-cycle when $\fsimp_i$ is positive.
    \item 
    $\Chnmat^{p-1}$ helps maintain the coboundary basis in $\Cobmat^{p}$.
    \item
     $\Rmat^{p-1}$ provides the boundary basis on which 
    pseudo-cochains in $\Chnmat^{p-1}$ can be defined.
\end{itemize}
% \end{remark}

% \Tao{The properties for these matrices need justification somewhere}
% (coming from the representatives of the ``active''
% harmonic bars, i.e., ending with $i$ as we sweep the filtration)

% \Tao{Need to mention that we represent the chains/cochains as $\Real$-``columns''}

\para{Initial reduction.}
% Our algorithm processes each insertion $K_{i}\inctosp{\fsimp_i}K_{i+1}$
% in $\filt$ one by one starting from $i=0$.
When processing the
insertion of 
a $p$-simplex $\fsimp_i$,
we first apply the  reduction  
% described 
in \Cref{algr:usual-reduce}
on $\Rmat^{p-1}$
(which is indeed 
% the one 
performed in the ordinary persistence algorithm~\cite{EdelsbrunnerLetscherZomorodian2002}).
% We have 
The following
\Cref{prop:usual-reduce} 
% which 
is  well-known~\cite{DeyWang2022,EdelsbrunnerLetscherZomorodian2002}:
% concerning the reduction:
\begin{proposition}
\label{prop:usual-reduce}
If 
% the reduction 
% on $\partial(\hat{\sG}_i)$ in 
Algorithm~\ref{algr:usual-reduce} ends with $z=0$,
then $\fsimp_i$ is positive;
otherwise, $\fsimp_i$ is negative.
\end{proposition}
Notice that besides determining whether $\fsimp_i$ is positive or negative,
% The purpose of 
% the reduction in 
\Cref{algr:usual-reduce}
% is two-fold:
% (1) 
% (2) 
also ensures that $\Rmat^{p-1}$ represents a basis for $\Bnd_{p-1}(K_{i+1})$
in the next iteration
(details discussed later).

\begin{algorithm}[h]
% {\sc Computing minimal persistent $\Dimless$-cycles of finite intervals for weak $(\Dimtop)$-pseudomanifolds}
\caption{Reduction based on $\Rmat^{p-1}$ for determining whether $\fsimp_i$ is positive or negative}

\begin{algorithmic}[1]
% \State $\phi\leftarrow\hat{\sG}_i$
\State $z\leftarrow\partial({\sG}_i)$
\While{$z\neq0$ and $\exists j$ s.t.\ $\pivot(\Rmat^{p-1}[j])=\pivot(z)$}
\State $k\leftarrow\pivot(z)$
\State $\alpha_1\leftarrow z[k]$
\State $\alpha_2\leftarrow\Rmat^{p-1}[j][k]$
\State $z\leftarrow z-(\alpha_1/\alpha_2)\cdot\Rmat^{p-1}[j]$
% \State $\phi\leftarrow \phi-(\alpha_1/\alpha_2)\cdot\Cobmat^{p}[j]$
\EndWhile
\end{algorithmic}
\label{algr:usual-reduce}
\end{algorithm}


% By \Cref{prop:reduce},
% % based on whether $\fsimp_i$ is positive or negative from the reduction,
% based on the result of the reduction,
% we follow different steps.

% \begin{itemize}
\para{$\sG_i$ is positive.}
We need to find a new harmonic $p$-cycle born in $K_{i+1}$.
For this, we perform another reduction on $\Bndmat^{p-1}$ and $\Cobmat^{p}$
 in \Cref{algr:reduce}.

\begin{algorithm}[!h]
% {\sc Computing minimal persistent $\Dimless$-cycles of finite intervals for weak $(\Dimtop)$-pseudomanifolds}
\caption{Reduction based on $\Bndmat^{p-1},\Cobmat^p$ for finding a new-born harmonic cycle 
% (repetition of \Cref{algr:mainbody-reduce} in \Cref{sec:algorithm})
}

\begin{algorithmic}[1]
\State $\phi\leftarrow\hat{\sG}_i$
\State $\zG\leftarrow\partial(\hat{\sG}_i)$
\While{$\zG\neq0$ and $\exists j$ s.t.\ $\pivot(\Bndmat^{p-1}[j])=\pivot(\zG)$}
\State $k\leftarrow\pivot(\zG)$
\State $\alpha_1\leftarrow\zG[k]$
\State $\alpha_2\leftarrow\Bndmat^{p-1}[j][k]$
\State $\zG\leftarrow \zG-(\alpha_1/\alpha_2)\cdot\Bndmat^{p-1}[j]$
\State $\phi\leftarrow \phi-(\alpha_1/\alpha_2)\cdot\Cobmat^{p}[j]$
\EndWhile
\end{algorithmic}
\label{algr:reduce}
\end{algorithm}

\begin{remark}
So far we have performed analogous operations on 
% the matrices 
$\Rmat^{p-1}$ and $\Bndmat^{p-1}$.
% One may  wonder whether one matrix can be replaced by the other
% especially considering that  both $\Rmat^{p-1}$, $\Bndmat^{p-1}$
% contain $(p-1)$-boundaries with distinct pivots.
However,
we notice that $\Rmat^{p-1}$ and $\Bndmat^{p-1}$
are in general two different matrices with different purposes.
The main purpose of $\Rmat^{p-1}$ is to form a filtered 
% but \emph{fixed}
basis for  $({p-1})$-boundary groups
of  complexes in $\filt$.
% With respect to these boundary basis, 
%  columns in $\Chnmat^{p-1}$ encode
% the $({p-1})$-pseudo-cochains.
% Hence, 
Columns in $\Rmat^{p-1}$ are {fixed}
once added to $\Rmat^{p-1}$.
On the other hand,
 $\Bndmat^{p-1}$ is mainly for 
finding a new-born harmonic cycle when $\sG_i$ is positive.
% {\red(by finding a new cocycle whose boundary is zero; see \Cref{algr:reduce})}.
Columns in
$\Bndmat^{p-1}$
are not fixed and can be changed in later steps
(such as in \Cref{algr:sum-bndmat}).
\end{remark}

\begin{proposition}
% [repetition of \Cref{prop:mainbody-reduce} in \Cref{sec:algorithm}]
\label{prop:reduce}
If 
% the reduction 
% on $\partial(\hat{\sG}_i)$ in 
Algorithm~\ref{algr:reduce} ends with $\zG=0$,
then $\fsimp_i$ is positive
and $\phi$ is a new harmonic $p$-cycle born in $\har_p(K_{i+1})$;
otherwise, $\fsimp_i$ is negative.
\end{proposition}
\begin{proof}

% First, 
Since $\hat{\fsimp}_i$ is a new $p$-cocycle in $K_{i+1}$,
% (see the proof of \Cref{prop:dim-chng})},
we have that $\phi$ is always a new $p$-cocycle in $K_{i+1}$
and  $\zG=\partial(\phi)$ at the end of each iteration in Algorithm~\ref{algr:reduce}.
% the reduction.
If Algorithm~\ref{algr:reduce}
% the reduction 
ends with $\zG=0$, then 
we have $\partial(\phi)=0$.
This implies that $\phi$ is a new harmonic $p$-cycle born in $\har_p(K_{i+1})$.
% and hence $\dim(\har_*(K_{i+1}))>\dim(\har_*(K_{i}))$.
% By definition, $\fsimp_i$ is positive.
By \Cref{thm:harmonic-inclusion},
$\fsimp_i$ is positive.

Now consider the case that \Cref{algr:reduce} ends with $\zG\neq 0$.
For  contradiction, suppose instead that $\fsimp_i$ is positive,
and let $\psi$ be a new harmonic $p$-cycle born in $\har_p(K_{i+1})$.
Since columns in $\Harmat^p$ and $\Cobmat^p$ form a basis of $\Zyc^p(K_i)$
and $\hat{\fsimp}_i$ is a new $p$-cocycle in $K_{i+1}$,
% {\red(see the proof of \Cref{prop:dim-chng})},
we have that $\Harmat^p$, $\Cobmat^p$, and $\hat{\fsimp}_i$ altogether form a basis of $\Zyc^p(K_{i+1})$.
Hence, 
% $\psi$ equals
% can be expressed as 
% a formal sum
we have
\mbeg\psi=\aG\cdot\hat{\fsimp}_i+\sum_j\mu_j\cdot\Harmat^p[j]+\sum_k\lG_k\cdot\Cobmat^p[k],\mend
where $\aG\neq 0$ because 
$\psi$ is a new  $p$-cocycle in $K_{i+1}$.
So we have \begin{align*}
0=\partial(\psi)&=\aG\cdot\partial(\hat{\fsimp}_i)+\sum_j\mu_j\cdot\partial(\Harmat^p[j])+\sum_k\lG_k\cdot\partial(\Cobmat^p[k])\\
&=\aG\cdot\partial(\hat{\fsimp}_i)+0+\sum_k\lG_k\cdot\Bndmat^{p-1}[k].
\end{align*}
Now \mbeg\partial(\hat{\fsimp}_i)=\sum_k(-\lG_k/\aG)\cdot\Bndmat^{p-1}[j],\mend
i.e., $\partial(\hat{\fsimp}_i)$ is a linear combination of
columns of $\Bndmat^{p-1}$.
Since the columns of $\Bndmat^{p-1}$ have distinct pivots,
the reduction in Algorithm~\ref{algr:reduce} must end with $\zG=0$,
which is a contradiction.
% because  there is not a column $\Bndmat^{p-1}[j]$ in 
\end{proof}

By \Cref{prop:reduce},
% since $\fsimp_i$ is positive,
% the reduction in 
\Cref{algr:reduce}
must end with $\zG=0$
and
% % Based on the proof of \Cref{prop:reduce},
% the cocycle $\phi$ at the end of the reduction 
% % in \Cref{algr:reduce}
% is a  new harmonic $p$-cocycle born in $K_{i+1}$.
% Accordingly,
we add $\phi$ as new column to $\Harmat^p$.
% \begin{proposition}
% % \Tao{state that each old coboundary is still a coboundary}
% Treating columns in $\Chnmat^{p-1}$
% and $\Cobmat^p$ as cochains in $K_{i+1}$,
% we still have that $\Cobmat^p[j]=\dG(\Chnmat^{p-1}[j])$
% for each $j$
% when the $p$-simplex $\fsimp_i$  is positive.
% \end{proposition}
% \begin{proof}
% For clarity,  let $\phi$ denote the coboundary of 
% $\Chnmat^{p-1}[j]$ in $K_{i+1}$.
% We only need to show that $\phi(\fsimp_i)=0$.
% Let $\{z_1,\ldots,z_\ell\}$ be a basis for $\Bnd_{p-1}(K_{i})$.
% Since $\fsimp_i$  is positive, we have that $\partial(\fsimp_i)\in\Bnd_{p-1}(K_{i})$
% (see, e.g., \cite{DeyWang2022,EdelsbrunnerLetscherZomorodian2002}),
% meaning that 
% $\partial(\fsimp_i)=\sum_{j=1}^\ell\aG_jz_j$,
% where each $\aG_j\in\Real$.
% We have
% $\phi(\fsimp_i)=\Chnmat^{p-1}[j](\partial(\fsimp_i))$.
% \end{proof}
% Notice that 
% when proceeding from $K_i$ to $K_{i+1}$,
% we need to {\red treat the pseudo-cochains encoded by $\Chnmat^{p-1}$
% and the cochains in $\Cobmat^p$ as ones defined in $K_{i+1}$}.
% \sout{\color{red} By proceeding from $K_i$ to $K_{i+1}$,
% cochains in $\Cobmat^p$ can be directly treated as cochains  in $K_{i+1}$
% due the identifications discussed before.}
% by {\red\bf the identification of cochains}.
Now
consider 
% an $\Chnmat^{p-1}[j]$ and 
the pseudo-cochain \mbeg\psi_j:\Bnd_{p-1}(K_i)\to\Real\mend in $K_i$
that a column $\Chnmat^{p-1}[j]$ encodes.
Since $\Bnd_{p-1}(K_i)=\Bnd_{p-1}(K_{i+1})$~\cite{DeyWang2022,EdelsbrunnerLetscherZomorodian2002},
% (a well-known fact~\cite{DeyWang2022,EdelsbrunnerLetscherZomorodian2002}),
we can directly treat $\psi_j$ as a pseudo-cochain   in $K_{i+1}$.
However, 
while we still have 
% $\dG(\Chnmat^{p-1}[j])(\fsimp_k)=\Cobmat^p[j](\fsimp_k)$
\mbeg\dG(\psi_j)(\fsimp)=\Cobmat^p[j](\fsimp)\text{ for each $p$-simplex }\fsimp\in K_i,\mend
we may have 
% that
\mbeg\dG(\psi_j)(\fsimp_i)=\psi_j(\partial\fsimp_i)\neq 0,\mend
% is non-zero,
% By doing this,
% $\dG(\Chnmat^{p-1}[j])(\fsimp_i)=\psi_j(\partial\fsimp_i)$
% for a certain $j$,
% by viewing 
so that \mbeg\Cobmat^p[j]\neq\dG(\Chnmat^{p-1}[j])\mend in $K_{i+1}$.
% (But notice that we still have 
% $\dG(\psi_j)(\fsimp_k)=\Cobmat^p[j](\fsimp_k)$
% for each $p$-simplex $\fsimp_k$
% with $k<i$
% because $\fsimp_k\in K_i$.)
% In this case, 
To deal with this,
we first allocate two new matrices 
$\bar{\Cobmat}^p$,
% and
$\bar{\Bndmat}^{p-1}$
% such that 
% making
and let
% $\bar{\Cobmat}^p[j]=\dG(\psi_j)$
\mbeg\bar{\Cobmat}^p[j]=\dG(\Chnmat^{p-1}[j]):=\dG(\psi_j)\text{, }\bar{\Bndmat}^{p-1}[j]=\partial(\bar{\Cobmat}^p[j])
\text{ for each $j$,}\mend
% in $K_{i+1}$
where 
% we let 
all the chains and (pseudo-)cochains
are now defined in $K_{i+1}$.
% where $\Chnmat^{p-1}[j]$ is now considered to 
% encode a pseudo-cochain
% where all the matrix columns are now viewed as chains or cochains in $K_{i+1}$.
To get $\bar{\Cobmat}^p[j]$,
% from ${\Cobmat}^p[j]$,
% as a coboundary of $\psi_j$,
we only need to know the value of $\psi_j(\partial\fsimp_i)$.
% Notice that 
From the reduction in \Cref{algr:usual-reduce},
we know how to express $\partial\fsimp_i$ as a formal sum of columns in $\Rmat^{p-1}$,
i.e., \mbeg\partial\fsimp_i=\sum_{k}\mu_k\Rmat^{p-1}[k],\mend
where each non-zero $\mu_k$ is given by the ``$\aG_1/\aG_2$'' in \Cref{algr:usual-reduce}.
Then,
\mbeg\psi_j(\partial\fsimp_i)=\psi_j(\sum_{k}\mu_k\Rmat^{p-1}[k])
=\sum_{k}\mu_k\psi_j(\Rmat^{p-1}[k])
=\sum_{k}\mu_k\Chnmat[j][k].\mend


Due to the previous update, columns in $\bar{\Bndmat}^{p-1}$ may not have distinct pivots.
We then 
% perform some updates 
run \Cref{algr:sum-bndmat}
% on  $\Chnmat^{p-1}$, $\bar{\Cobmat}^p$, and $\bar{\Bndmat}^{p-1}$
to make columns in $\bar{\Bndmat}^{p-1}$ have distinct pivots again.
% After this, we let 
% ${\Cobmat}^p=\bar{\Cobmat}^p$
% and
% ${\Bndmat}^{p-1}=\bar{\Bndmat}^{p-1}$.
% In this case, we cannot simply update 
% $\Cobmat^p[j]$ as $\dG(\Chnmat^{p-1}[j])$
% and accordingly update
% $\Bndmat^{p-1}[j]$ as $\partial(\Cobmat^p[j])$
% for each such $j$,
% because pivots in the columns $\Bndmat^{p-1}[j]$ may not be distinct
% after the update.
% % We further notice
% % it would be too costly to perform additional updates to make 
% % the pivot distinct again because there may be more than two columns
% % in $\Bndmat^{p-1}[j]$ .
To see how \Cref{algr:sum-bndmat} works,
% let ${\Cobmat}^p[\lG_j]=\sum_{k=0}^{i-1}\gG_k\hat{\fsimp}_k$
% for $j\in\{2,\ldots,\ell\}$.
% Then $\bar{\Cobmat}^p[\lG_j]=\sum_{k=0}^{i-1}\gG_k\hat{\fsimp}_k+\aG_j\hat{\fsimp_i}$,
notice that before executing the for loop in line~\ref{line:sum-bndmat-for}, we have
% \Cref{algr:sum-bndmat},
\mbeg\bar{\Cobmat}^p[\lG_j]={\Cobmat}^p[\lG_j]+\aG_j\hat{\fsimp}_i
\text{ for }j\in\{1,\ldots,\ell\}.\mend
This means that
\mbeg\bar{\Bndmat}^{p-1}[\lG_j]={\Bndmat}^{p-1}[\lG_j]+\aG_j\partial(\hat{\fsimp}_i).\mend
So for $j>1$, we have
\begin{align*}
&\bar{\Bndmat}^{p-1}[\lG_j]-(\aG_j/\aG_1)\cdot\bar{\Bndmat}^{p-1}[\lG_1]\\
=&{\Bndmat}^{p-1}[\lG_j]+\aG_j\partial(\hat{\fsimp}_i)
-(\aG_j/\aG_1)({\Bndmat}^{p-1}[\lG_1]+\aG_1\partial(\hat{\fsimp}_i))\\
=&{\Bndmat}^{p-1}[\lG_j]-(\aG_j/\aG_1)\cdot{\Bndmat}^{p-1}[\lG_1].
\end{align*}
Since $\pivot({\Bndmat}^{p-1}[\lG_1])<\pivot({\Bndmat}^{p-1}[\lG_j])$,
we have \mbeg\pivot(\bar{\Bndmat}^{p-1}[\lG_j])=\pivot({\Bndmat}^{p-1}[\lG_j])\mend
after executing the for loop.
% in line~\ref{line:sum-bndmat-for}.
Now we have \mbeg\bar{\Bndmat}^{p-1}[\lG]={\Bndmat}^{p-1}[\lG]\text{ for }\lG\not\in\{\lG_1,\ldots,\lG_\ell\},\mend
meaning that
before executing the while loop
in line~\ref{line:sum-bndmat-while},
only pivots of the two columns indexed by $\lG_1$ may differ in
$\bar{\Bndmat}^{p-1}$ and ${\Bndmat}^{p-1}$.
This means that  at most two columns in
$\bar{\Bndmat}^{p-1}$ may have the same pivot
before executing the while loop.
% Since each iteration of the loop changes only one column of
% $\bar{\Bndmat}$,
We then have that,
in each iteration of the while loop,
at most two columns in
$\bar{\Bndmat}^{p-1}$ have the same pivot.
Hence, columns in $\bar{\Bndmat}^{p-1}$
have distinct pivots after finishing the loop.
Notice that the pivot $k$ in line~\ref{line:sum-bndmat-k}
decreases in each iteration,
meaning that  the while loop executes for no more than $i$ iterations.
Hence, \Cref{algr:sum-bndmat} runs in $O(\filtcnt^2)$ time.
Notice that allocating the two new matrices
$\bar{\Cobmat}^p$,
 $\bar{\Bndmat}^{p-1}$
are purely for clearly presenting and justifying the operations
in \Cref{algr:sum-bndmat}.
To implement this in computer, 
the operations can be directly performed  on 
${\Cobmat}^p$
and ${\Bndmat}^{p-1}$.

\begin{algorithm}[h]
% {\sc Computing minimal persistent $\Dimless$-cycles of finite intervals for weak $(\Dimtop)$-pseudomanifolds}
\caption{Making columns in $\bar{\Bndmat}^{p-1}$ have distinct pivots}

\begin{algorithmic}[1]
\State let $\bar{\Cobmat}^p[\lG_1],\ldots,\bar{\Cobmat}^p[\lG_\ell]$ be all the columns in
$\bar{\Cobmat}^p$ s.t.\ $\bar{\Cobmat}^p[\lG_j](\fsimp_i)\neq 0$
for each $j$
\State reorder $\lG_1,\ldots,\lG_\ell$ s.t.\ $\pivot({\Bndmat}^{p-1}[\lG_1])<\pivot({\Bndmat}^{p-1}[\lG_j])$
for each $j>1$
\State $\aG_1\leftarrow\bar{\Cobmat}^p[\lG_1](\fsimp_i)$
\For{$j=2,\ldots,\ell$}\label{line:sum-bndmat-for}
\State $\aG_j\leftarrow\bar{\Cobmat}^p[\lG_j](\fsimp_i)$
\State ${\Chnmat}^{p-1}[\lG_j]\leftarrow{\Chnmat}^{p-1}[\lG_j]-(\aG_j/\aG_1)\cdot{\Chnmat}^{p-1}[\lG_1]$
\State $\bar{\Cobmat}^p[\lG_j]\leftarrow\bar{\Cobmat}^p[\lG_j]-(\aG_j/\aG_1)\cdot\bar{\Cobmat}^p[\lG_1]$
\State $\bar{\Bndmat}^{p-1}[\lG_j]\leftarrow\bar{\Bndmat}^{p-1}[\lG_j]-(\aG_j/\aG_1)\cdot\bar{\Bndmat}^{p-1}[\lG_1]$
\EndFor
% \LineComment{\textsf{\textit{Find any two columns with same pivot in $\bar{\Bndmat}^{p-1}$}}}
\While{exist two columns $\bar{\Bndmat}^{p-1}[\lG]$, $\bar{\Bndmat}^{p-1}[\mG]$ with same pivot}\label{line:sum-bndmat-while}
% \While{$\exists \lG,\mG$ s.t.\ $\pivot(\bar{\Bndmat}^{p-1}[\lG])=\pivot(\bar{\Bndmat}^{p-1}[\mG])$}\label{line:sum-bndmat-while}
\State $k\leftarrow\pivot(\Bndmat^{p-1}[\lG])$\label{line:sum-bndmat-k}
\State $\gamma_1\leftarrow\Bndmat^{p-1}[\lG][k]$
\State $\gamma_2\leftarrow\Bndmat^{p-1}[\mG][k]$
\State ${\Chnmat}^{p-1}[\lG]\leftarrow{\Chnmat}^{p-1}[\lG]-(\gamma_1/\gamma_2)\cdot{\Chnmat}^{p-1}[\mG]$
\State $\bar{\Cobmat}^p[\lG]\leftarrow\bar{\Cobmat}^p[\lG]-(\gamma_1/\gamma_2)\cdot\bar{\Cobmat}^p[\mG]$
\State $\bar{\Bndmat}^{p-1}[\lG]\leftarrow\bar{\Bndmat}^{p-1}[\lG]-(\gamma_1/\gamma_2)\cdot\bar{\Bndmat}^{p-1}[\mG]$
\EndWhile
\State{${\Cobmat}^p\leftarrow\bar{\Cobmat}^p$}
\State{${\Bndmat}^{p-1}\leftarrow\bar{\Bndmat}^{p-1}$}
\end{algorithmic}
\label{algr:sum-bndmat}
\end{algorithm}

Finally, we have that
the other matrices do not need to be changed
% which follows from 
by \Cref{prop:dim-chng}.
% This follows from \Cref{thm:harmonic-inclusion}, \Cref{prop:dim-chng},
% and some well-know facts concerning ordinary persistence~\cite{DeyWang2022,EdelsbrunnerLetscherZomorodian2002}.

% \Tao{Argue that the coboudary matrix is still a basis? This may be 'trivial' but probably needs a justification?}

% \Tao{may need to rephrase the proposition})
% Since
% $\har_q(K_{i+1})=\har_q(K_{i})$ for $q\neq p$
% and 
% $\Bnd^q(K_{i+1})=\Bnd^q(K_{i})$
% for each $q$ (see \Cref{prop:dim-chng}),
% the other matrices do not need to be changed.
% \Tao{we need to update the p-coboundary matrix}

\para{$\sG_i$ is negative.}
First update $\Harmat^{p-1}$ 
according to \Cref{alg:frame}.
% Since columns in $\Harmat^{p-1}$ are tentative representatives for the unpaired birth indices,
% we remove the {\red column corresponding to $b_{j*}^{p-1}$ 
% from $\Harmat^{p-1}$} 
% For this, 
Since $\fsimp_i$ is negative, 
we have $z\neq 0$ at the end of \Cref{algr:usual-reduce}.
We then add $z$ to $\Rmat^{p-1}$ 
to form a basis for $\Bnd_{p-1}(K_{i+1})$
(see~\cite{DeyWang2022,EdelsbrunnerLetscherZomorodian2002}).
For ease of presentation, 
let $\Rmat^{p-1}$ still denote 
% the version of 
the matrix
before $z$ is added
and let $\bar{\Rmat}^{p-1}=\Rmat^{p-1}\cup\{z\}$.
We extend each column 
% $\Chnmat^{p-1}[j]$ 
% $\Chnmat^{p-1}$ 
of $\Chnmat^{p-1}$
by appending 0
so that columns of $\Chnmat^{p-1}$
encode $(p-1)$-pseudo-cochains in $K_{i+1}$
w.r.t\ $\bar{\Rmat}^{p-1}$.
% This means that $\dG(\psi_j)(z)=0$
% where $\psi_j$
% % $\psi_j:\Bnd_{p-1}(K_{i+1})\to\Real$ 
% is the 
% % $(p-1)$-
% pseudo-cochain
% encoded by  $\Chnmat^{p-1}[j]$ w.r.t $\bar{\Rmat}^{p-1}$.
Notice that we may also have
$\dG(\Chnmat^{p-1}[j])(\fsimp_i)\neq 0$
for a $j$.
% $\dG(\psi_j)(\fsimp_i)=\psi_j(\partial\fsimp_i)\neq 0$.
% Similar to the previous case where $\fsimp_i$ is positive,
% the reduction in \Cref{algr:usual-reduce} also provides
% a way to express $\partial\fsimp_i$
% as a formal sum of columns in $\bar{\Rmat}^{p-1}$.
Hence,
we update 
$\Cobmat^p$ and $\Bndmat^{p-1}$
as in the previous case where $\fsimp_i$ is positive,
 % where $\fsimp_i$ is positive
including performing \Cref{algr:sum-bndmat}.
After this, 
${\Cobmat}^p[j]=\dG(\Chnmat^{p-1}[j])$, ${\Bndmat}^{p-1}[j]=\partial({\Cobmat}^p[j])$
in $K_{i+1}$
% in $K_{i+1}$
for each $j$
and columns in ${\Bndmat}^{p-1}$ have distinct pivots.

Since $\dim(\Bnd^{p}(K_{i+1}))=\dim(\Bnd^{p}(K_{i}))+1$
by \Cref{prop:dim-chng},
% and so
we need to add a new column to $\Cobmat^p$ 
% (and accordingly to $\Bndmat^{p-1}$)
% so that $\Cobmat^p$ 
to form a basis for $\Bnd^{p}(K_{i+1})$.
For this,
we first add the column 
\mbeg\vec{v}=(0,\ldots,0,1)^T\mend
% $\vec{v}=(\eta_1=0,\ldots,\eta_{r-1}=0,\eta_{r}=1)^T$
to $\Chnmat^{p-1}$
and observe the following:
% where $r$ is the number of columns in $\bar{\Rmat}^{p-1}$.
\begin{proposition}\label{prop:sG-i-new-cob}
Let $\vec{v}$ be as defined above.
Then,
$\hat{\fsimp}_i=\dG(\vec{v})$.
\end{proposition}
\begin{proof}
Let $r$ be the length of $\vec{v}$
and let $\psi$ be the pseudo-cochain encoded by $\vec{v}$
w.r.t 
% the basis 
$\bar{\Rmat}^{p-1}$.
For a $p$-simplex $\fsimp\in K_i$,
we have 
\mbeg
\dG(\vec{v})(\fsimp)=\psi(\partial\fsimp)=\psi(\sum_{k=1}^{r-1}\mu_k{\Rmat}^{p-1}[k])
=\sum_{k=1}^{r-1}\mu_k\psi(\Rmat^{p-1}[k])
% =\sum_{k=1}^{r-1}\mu_k\eta_k
=0.
\mend
Also, by the reduction in \Cref{algr:usual-reduce},
we have \mbeg\partial\fsimp_i=\sum_{k=1}^{r-1}\lG_k{\Rmat}^{p-1}[k]+z.\mend
Then,
\mbeg\dG(\vec{v})(\fsimp_i)=\psi(\partial\fsimp_i)=\psi(\sum_{k=1}^{r-1}\lG_k{\Rmat}^{p-1}[k]+z)
=\psi(z)=1.\qedhere\mend
\end{proof}

\begin{proposition}\label{prop:add-to-cob-mat}
% {\red After executing \Cref{algr:sum-bndmat} when $\fsimp_i$ is negative,}
The columns in ${\Cobmat}^p$ 
% be the matrix
% resulting from adding the column 
along with 
% the $p$-cochain 
$\hat{\fsimp}_i$ 
% to ${\Cobmat}^p$.
% One has that $\bar{\Cobmat}^p$
form a basis for $\Bnd^p(K_{i+1})$.
\end{proposition}
\begin{proof}
% Let
Abuse the notations slightly by
letting
${\Cobmat}^p$
and $\bar{\Cobmat}^p$ denote  versions of the two matrices before executing \Cref{algr:sum-bndmat},
i.e., columns in ${\Cobmat}^p$ are still cochains in $K_i$ 
forming a basis of $\Bnd^p(K_i)$
and
each $\bar{\Cobmat}^p[j]$ (a cochain in $K_{i+1}$)
% of  $\bar{\Cobmat}^p$ 
is simply derived by taking the coboundary 
of $\Chnmat^{p-1}[j]$ in  $K_{i+1}$.
We first show that columns of
$\bar{\Cobmat}^p[j]$ and $\hat{\fsimp}_i$  are linearly 
independent.
Suppose that \mbeg\sum_j\aG_j\bar{\Cobmat}^p[j]+\gamma\hat{\fsimp}_i=0.\mend
Let $\iG^\sharp:\Chn^p(K_{i+1})\to\Chn^p(K_{i})$
be the cochain map induced by inclusion.
% We have $\iG^\sharp(\bar{\Cobmat}^p[j])={\Cobmat}^p[j]$
% for each $j$.
We have that \mbeg\sum_j\aG_j{\Cobmat}^p[j]=\sum_j\aG_j\iG^\sharp(\bar{\Cobmat}^p[j])
=\iG^\sharp(\sum_j\aG_j\bar{\Cobmat}^p[j])
=\iG^\sharp(-\gamma\hat{\fsimp}_i)
=-\gamma\iG^\sharp(\hat{\fsimp}_i)=0,\mend
meaning that each $\aG_j$ is zero.
This also means that $\gamma=0$  because $\hat{\fsimp}_i\neq 0$.
We then have that
columns of
$\bar{\Cobmat}^p[j]$ along with $\hat{\fsimp}_i$
form a basis of $\Bnd^{p}(K_{i+1})$
because $\hat{\fsimp}_i\in\Bnd^{p}(K_{i+1})$
(\Cref{prop:sG-i-new-cob}) and $\dim(\Bnd^{p}(K_{i+1}))=\dim(\Bnd^{p}(K_{i}))+1$ (\Cref{prop:dim-chng}).
% The proposition 
\Cref{prop:add-to-cob-mat}
then follows because   \Cref{algr:sum-bndmat}
does not change the linear span of matrix columns.
\end{proof}

By \Cref{prop:sG-i-new-cob,prop:add-to-cob-mat},
we add $\hat{\fsimp}_i$
to $\Cobmat^p$
and add $\partial(\hat{\fsimp}_i)$
to $\Bndmat^{p-1}$.
Notice that there may be at most two columns having the same pivot
in $\Bndmat^{p-1}$ after this.
% adding $\partial(\hat{\fsimp}_i)$.
For this, we apply a procedure similar to the while loop in 
\Cref{algr:sum-bndmat}
to make $\Bndmat^{p-1}$ have distinct pivots again.

By \Cref{prop:dim-chng},
% We also that
% Notice that
the other matrices do not need to be changed.
% (see \Cref{thm:harmonic-inclusion}, {\red\Cref{prop:dim-chng}}, and~\cite{DeyWang2022,EdelsbrunnerLetscherZomorodian2002}).

\bigskip

We now  conclude the following:
% the section:
\begin{theorem}
% [repetition of \Cref{thm:mainbody-m3} in \Cref{sec:algorithm}]
    The \hcb{} of $\filt$ can be computed in $O(\filtcnt^3)$ time.
\end{theorem}
\begin{proof}
Each iteration of the algorithm takes no more than $O(\filtcnt^2)$ time
and so
the overall complexity 
% of the algorithm 
is $O(\filtcnt^3)$.
Correctness of our implementation 
% presented in \Cref{set:impl} 
follows from 
correctness of the operations detailed in this section.
% \Cref{sec:impl}.
% and the fact that properties of the matrices
% claimed at the beginning of \Cref{sec:impl}
% are correctly maintained.
\end{proof}



% Doing this relies on the following fact:

% \begin{proposition}\label{prop:sG-i-new-cob}
% If  the $p$-simplex $\fsimp_i$ 
% % added  from $K_i$ to $K_{i+1}$
% is negative, then $\hat{\fsimp}_i$ is a new $p$-coboundary in $K_{i+1}$.
% \end{proposition}
% \begin{proof}
% Since $\hat{\fsimp}_i$ is  a new cocycle in $K_{i+1}$,
% we only need to show that $\hat{\fsimp}_i$ is a coboundary in $K_{i+1}$.
% Let $\{z_1,\ldots,z_\ell\}$ be a basis for $\Bnd_{p-1}(K_{i})$.
% Since $\fsimp_i$  is negative, we have that $\partial(\fsimp_i)\not\in\Bnd_{p-1}(K_{i})$
% (see, e.g., \cite{DeyWang2022,EdelsbrunnerLetscherZomorodian2002}).
% Hence, $\Bcal=\{z_1,\ldots,z_\ell,\partial(\fsimp_i)\}$ forms a basis for $\Bnd_{p-1}(K_{i+1})$.
% We can then extend $\Bcal$ into a basis $\Bcal'$
% for $\Chn_{p-1}(K_{i+1})$ by adding additional chains.
% Define a cochain $\phi\in\Chn^{p-1}(K_{i+1})$
% where $\phi(\partial(\fsimp_i))=1$ and  $\phi(z)=0$ for each $z\in\Bcal'\setminus\{\partial(\fsimp_i)\}$.
% We have $\delta(\phi)(\fsimp_i)=\phi(\partial(\fsimp_i))=1$.
% Also, for each $p$-simplex $\sG\neq\fsimp_i$ in $K_{i+1}$,
% since $\partial(\sG)\in\Bnd_{p-1}(K_i)$,
% we have $\delta(\phi)(\sG)=\phi(\partial(\sG))=\phi(\sum_{j=1}^\ell\aG_jz_j)=0$,
% where each $\aG_j\in\Real$.
% Hence, $\hat{\fsimp}_i=\delta(\phi)$
% and so $\hat{\fsimp}_i$ is a coboundary in $K_{i+1}$.
% \end{proof}

% Notice that by proceeding from $K_i$ to $K_{i+1}$,
% we may also have
% $\Cobmat^p[j]\neq\dG(\Chnmat^{p-1}[j])$ for a $j$ in $K_{i+1}$ as in the previous case. 
% % that $\fsimp_i$ is positive.
% Hence, we also apply \Cref{algr:sum-bndmat} 
% so that
% $\Cobmat^p[j]=\dG(\Chnmat^{p-1}[j])$ for each $j$ 
% and columns in $\Bndmat^{p-1}$ have distinct pivots.
% Let $\fsimp_\lG$ be a $(p-1)$-face of $\fsimp_i$
% (notice that $\lG<i$).
% We add a new column $\zG$ to 
% $\Chnmat^{p-1}$
% where $\zG[\lG]=1$
% and $\zG[k]=0$ for $k\neq \lG$.
% Accordingly, we add $\dG(\zG)$ to $\Cobmat^p$
% and $\partial\circ\dG(\zG)$ to $\Bndmat^{p-1}$.
% Notice that since $\dG(\zG)(\fsimp_i)=1$,

% \begin{proposition}\label{prop:pers-pair-cob}
% Suppose that the $p$-simplex $\fsimp_i$ 
% is negative.
% Let $(\fsimp_j,\fsimp_i)$
% be the ``persistence pair'' containing $\fsimp_i$ as
% in the usual persistent homology~\cite{EdelsbrunnerLetscherZomorodian2002}.
% One has $\hat{\fsimp}_i=\dG(\hat{\fsimp}_j)$,
% where $\hat{\fsimp}_i$ and $\hat{\fsimp}_j$ are both
% cochains in $K_{i+1}$.
% \end{proposition}
% \begin{proof}
% Computing  the usual persistent homology boils
% down to computing an $R=DV$ decomposition~\cite{cohen2006vines},
% where $D$ is the persistence boundary matrix 
% and $R$ is a ``reduced'' matrix
%  containing a basis for 
%  the boundary group
%  % $\Bnd_*(K_k)$
% % with distinct pivots 
% in each iteration
% of the algorithm.
% Specifically, let $R$ be the reduced matrix 
% after processing $\fsimp_i$ 
% when running the persistence algorithm~\cite{cohen2006vines,EdelsbrunnerLetscherZomorodian2002}
% on $\filt$.
% Since columns in $R$ correspond to simplices in $\filt$,
% we index $R$'s columns 
% % in $R$ 
% by their corresponding simplices,
% % in $\filt$,
% i.e.,
% $R[\fsimp_k]$ denotes the column 
% % of $R$
% corresponding to 
% % the simplex 
% $\fsimp_k$.
% % in $\filt$.
% For simplicity, assume $R$ contains only 
% columns corresponding to  
% % $p$-simplices.
% $\fsimp_i$ and those $p$-simplices added before $\fsimp_i$
% in $\filt$.
% We have that all non-zero columns in $R$ (which have distinct pivots) form a basis of $\Bnd_{p-1}(K_{i+1})$,
% and non-zero columns in $R$ other than $R[\fsimp_i]$ form a basis of $\Bnd_{p-1}(K_{i})$.
% Let $\pi=\pivot(R[\fsimp_i])$,
% and notice that the $\pi$-th
% row of $R$ corresponds to the $(p-1)$-simplex $\fsimp_j$.
% For each non-zero column $R[\fsimp_k]$ such that $\pivot(R[\fsimp_k])>\pi$
% and $R[\fsimp_k][\pi]\neq 0$,
% add $\aG\cdot R[\fsimp_i]$ 
% % (times a coefficient)
% to $R[\fsimp_k]$ to make
% $R[\fsimp_k][\pi]= 0$,
% where $\aG$ is a properly chosen coefficient.
% By doing this,
% % for each such $R[\fsimp_k]$,
% we get a reduced matrix $R$ where the $\pi$-th row
% is zero everywhere except in the column $R[\fsimp_i]$.
% We then have 
% $\hat{\fsimp}_j(R[\fsimp_i])=1$
% and
% $\hat{\fsimp}_j(R[\fsimp_k])=0$ for 
% % each $R[\fsimp_k]$ with 
% $k\neq i$.
% Notice that the processing of $\fsimp_i$ in the persistence algorithm~\cite{cohen2006vines,EdelsbrunnerLetscherZomorodian2002}
% indeed runs a reduction (similar to \Cref{algr:reduce}) 
% on 
% % $\partial(\fsimp_i)$ and 
% columns in $R$
% with $R[\fsimp_i]=\partial(\fsimp_i)$ at the beginning.
% The reduction keeps
% adding other columns of $R$ to $R[\fsimp_i]$
% until the pivot of $R[\fsimp_i]$
% does not equal the pivot of any other column of $R$.
% Then we have that $\partial(\fsimp_i)=\sum_{k}\aG_k\cdot R[\fsimp_k]$
% where $\aG_i=1$,
% which implies that 
% $\hat{\fsimp}_j(\partial(\fsimp_i))=\sum_{k}\aG_k\cdot \hat{\fsimp}_j(R[\fsimp_k])=1$.
% Since non-zero columns in $R$ other than $R[\fsimp_i]$ form a basis of $\Bnd_{p-1}(K_{i})$,
% we have that for $\sG\neq\fsimp_i\in K_{i+1}$,
% $\partial(\sG)=\sum_{k\neq i}\gG_k\cdot R[\fsimp_k]$
% because $\partial(\sG)\in\Bnd_{p-1}(K_{i})$.
% So we have 
% $\hat{\fsimp}_j(\partial(\sG))=\sum_{k\neq i}\gG_k\cdot \hat{\fsimp}_j(R[\fsimp_k])=0$.
% Hence, $\hat{\fsimp}_i=\dG(\hat{\fsimp}_j)$.
% \end{proof}

% By \Cref{prop:pers-pair-cob,prop:add-to-cob-mat},
% we add the column $\hat{\fsimp}_i$ to  $\Cobmat^p$,
% the column $\hat{\fsimp}_j$ to  $\Chnmat^{p-1}$,
% and the column  $\partial(\hat{\fsimp}_i)$ to  $\Bndmat^{p-1}$.




% By \Cref{prop:sG-i-new-cob}, since $\hat{\fsimp}_i$ is negative, 
% $\phi$ is initially a new $p$-coboundary in $K_{i+1}$ in the reduction in \Cref{algr:reduce}.
% Since we only add coboundaries to $\phi$ in each iteration,
% % of the reduction,
% $\phi$ is still a new $p$-coboundary in $K_{i+1}$
% at the end of the reduction.
% END AUTHOR FILE sec-barcode-zigzag-algo-impl.tex
% END AUTHOR FILE sec-barcode-zigzag-algo.tex
% \input{sec-barcode-zigzag-algo-impl}
% BEGIN AUTHOR FILE sec-stability.tex
\section{Sublevel Set Harmonic Chain Barcode and Its Stability}
\label{sec:stability}

In this section, we introduce the notion of \emph{sublevel set harmonic chain barcodes} and present our stability results based on this notion.
% ; see full details in \Cref{sec:stability}. 
Our stability proof makes use of the work on the algebraic stability of block decomposable $\Rspace^2$-modules~\cite{bjerkevik2021stability,botnan2018algebraic}. 
In brief, we ``lift'' the 1-parameter harmonic zigzag module to an $\Rspace^2$-indexed 2-parameter persistence module which is block decomposable. 
We then show that in the typical setting of sublevel set filtrations, there exists an interleaving between the lifted modules. Our main work here is to construct a concrete extension (see~\cref{subsec:extension}) and an interleaving between extensions to $\Rspace^2$-modules (see~\cref{subsec:stability}) realizing the category-theoretic concepts employed in~\cite{bjerkevik2021stability,botnan2018algebraic}. From the Isometry Theorem~\cite{bauer2014induced,botnan2018algebraic}, we then deduce the stability of the sublevel set {\hcb}.
We also address the disparity between the natural {\hcb} of a sublevel set filtration (which consists of closed-open intervals\footnote{The closed-open intervals over the real values we have here are different from the closed-open intervals of~\cite{bjerkevik2021stability,botnan2018algebraic} over the integers. Indeed, when considering only integer indices, we only have closed-closed intervals in our setting. The extension of our closed-open interval would be similar to a closed block of~\cite{botnan2018algebraic} but with the right vertical side open.}) and common conventions of zigzag persistence (which work with closed-closed intervals in our setting).  

\subsection{Interleaving and Stability for Ordinary Persistence} \label{subsec:ordinaryinterleaving}


Let $F$ and $G$ be two filtrations over the complexes $K$ and $K'$ respectively, and let the maps connecting the complexes in $F$ and $G$ be $f^{s,t}$ and $g^{s,t}$ respectively, for  $s\leq t \in \Rspace$. Let $C(F)$ and $C(G)$ be the corresponding filtrations of chain groups, and $H(F)$ and $H(G)$ the corresponding persistence modules. With an abuse of notation, we denote the maps induced on chain groups and homology groups also by $f^{s,t}$ and $g^{s,t}$ respectively, where $s\leq t \in \Rspace$.

\begin{definition}
\label{def:chain-interleaving}
Let $F$ and $G$ be two filtrations over  $K$ and $K'$ respectively. 
An \emph{$\eps$-chain-interleaving between $F$ and $G$} (or an $\eps$-interleaving at the chain level) is given by two sets of homomorphisms $\{\phi_\alpha : C(K_\alpha) \xrightarrow{} C(K'_{\alpha+\eps}) \}$ and $\{\psi_\alpha : C(K'_{\alpha}) \xrightarrow{} C(K_{\alpha+\eps}) \}$, such that 
\begin{enumerate}
    \item $\{ \phi_\alpha \}$ and $\{\psi_\alpha \}$ commute with the maps in the filtration, that is, for all $\alpha,t \in \Rspace$, $g^{\alpha+\eps,\alpha+\eps + t} \phi_\alpha = \phi_{\alpha+t} f^{\alpha, \alpha+t}$ and $f^{\alpha+\eps,\alpha+\eps+t} \psi_\alpha = \psi_{\alpha+t} g^{\alpha, \alpha+t}$; 
    \item The following diagrams commute: 
    \begin{equation}
      \begin{tikzcd}
  C(K_\alpha)
    \ar[r,"f^{\alpha, \alpha+\eps}"]
    \ar[dr, "\phi_\alpha", very near start, outer sep = -2pt]
  & C(K_{\alpha+\eps})
    \ar[r, "f^{\alpha+\eps, \alpha+2\eps}"]
    \ar[dr, "\phi_{\alpha+\eps}", very near start, outer sep = -2pt]
  & C(K_{\alpha+2\eps})
  \\
  C(K'_\alpha)
    \ar[r, swap,"g^{\alpha, \alpha+\eps}"]
    \ar[ur, crossing over, "\psi_\alpha"', very near start, outer sep = -2pt]
  & C(K'_{\alpha+\eps})
    \ar[r, swap,"g^{\alpha+\eps, \alpha+2\eps}"]
    \ar[ur, crossing over, "\psi_{\alpha+\eps}"', very near start, outer sep = -2pt]
  & C(K'_{\alpha+2\eps}).
  \end{tikzcd}
    \end{equation}
\end{enumerate}
\noindent The \emph{chain interleaving distance} is defined to be 
\begin{equation}
d_{CI}(F, G) := \inf \{\eps\geq 0 \mid \text{there exists an $\eps$-chain interleaving between } F\; \text{and}\; G\}.
\end{equation}
\end{definition}

We notice that the standard notion of $\eps$-interleaving~\cite{ChazalCohenSteinerGlisse2009}, denoted $d_I(F,G)$, is analogous to our definition above. However, it is defined on the persistence modules $H(F)$ and $H(G)$ on the homology level~\cite{ChazalCohenSteinerGlisse2009}, rather than the filtration of chain groups.  
Filtrations on the chain level arising from sublevel-sets of a function are ones that interest us in this paper.
% The filtration on the chain level of chains exists in the sublevel-set application that interests us.

% The stability of persistence diagrams (or barcodes) is a crucial property for applications. It says that small changes in data lead to small changes in the persistence diagrams (PDs). While the interleaving distance provides a measure of change in data, we measure the distance between PDs using the bottleneck distance. Let $D$ and $D'$ denote two PDs. 
% Recall that a PD is a multi-set of points in the extended plane (each of which is a birth-death pair) which also contains all points on the diagonal.
% The \emph{bottleneck distance} is defined as 
% \[
% d_B(D, D') = \text{inf}_\gamma \text{sup}_{p \in D} || p - \gamma(p) ||_\infty,
% \]
% where $\gamma$ ranges over all bijections between $D$ and $D'$ and $||\cdot ||_\infty$ is the largest absolute value of differences of the points' coordinates.
% % (or  the largest distortion of an endpoint 
% % % when matched to the bar of $D'$ 
% % when considering the PD points as bars).   

We refer to~\cite{ChazalCohenSteinerGlisse2009} for proof of the following 
stability theorem for ordinary persistence.
% \Cref{theorem:persistence-stability}.

\begin{theorem}
\label{theorem:persistence-stability}
    Let $F$ and $G$ be filtrations defined over (finite) complexes $K$ and $K'$, respectively. Then 
    $$ d_B( \Dgm(F), \Dgm(G)) \leq d_I(F,G).$$
\end{theorem}

See~\cite{Cohen-SteinerEdelsbrunnerHarer2005, ChazalCohenSteinerGlisse2009} for more on the stability of ordinary persistence. We also notice that in the above theorem we can replace $d_{CI}$ in place of $d_I$, since the existence of an $\eps$-chain-interleaving implies the existence of an $\eps$-interleaving on the homology level. 
% The resulting theorem would be weaker though.


\subsection{Sub-level-Set Filtration}\label{sec:Sub-level-Set-Filtration}

Note that we assumed up until now a filtration to be simplex-wise. While this is not true for a sublevel set filtration, we shall construct simplex-wise filtrations from a general filtration. 
% \sout{\red From now on, $Z_p(K)$ and $Z_{*}(K)$ denotes the group of cycles of $K$ for a fixed dimension $p$ and for all dimensions, respectively.}


\begin{definition}\label{dfn:bd-span-bar2}
Let $\tilde{F}$ be an increasing filtration of complex $K$ which is not necessarily simplex-wise and let $z \in Z_p(K)$ be a $p$-cycle. 
Let $t_{b(z)}$ be the time when $z$ is born,  i.e., $z\in Z_p(K_{t_{b(z)}})\setminus Z_p(K_{t_{b(z)-1}})$, and $t_{d(z)}$ be the time when $z$ dies as a harmonic chain, i.e., $\delta(z)\neq 0 $ in $K_{t_{d(z)}}$, but $\delta(z)=0$ in $K_t$ for $t_{b(z)} \leq t < t_{d(z)}$. Define the \emph{span} of $z$ as $span(z)=[t_{b(z)}, t_{d(z)})$.  If $z$ is not harmonic at birth, its span is the empty set (notice that a cycle  continues to be non-harmonic once it becomes so).
    The \emph{bar} of $z$ is defined as $bar(z) := [t_{b(z)}, t_{d(z)-1}]$.
\end{definition}


Let $\tilde{F}$ be an increasing filtration of $K$ which is not necessarily simplex-wise, i.e., two consecutive complexes in $\tilde{F}$ can differ by more than one simplex. We fix an ordering, once and for all, for vertices of $K$, which also fixes an ordering for all $K$'s simplices (using, say, the lexicographic ordering of the ordered sequence of vertices in a simplex). Given $\delta>0$ small enough, 
we expand $\tilde{F}$ into a simplex-wise filtration $\tilde{F}(\delta)$
by  expanding each inclusion $K_{t_{i-1}}\incto K_{t_{i}}$
in $\tilde{F}$
starting with $i=1$.
To differentiate,
 denote each complex in $\tilde{F}(\delta)$
as $K'_t$ where $t\in\mathbb{R}$ is a timestamp.
In each iteration, we expand the following inclusion in $\tilde{F}$:
$$K_{t_{i-1}} \incto K_{t_i} 
% \xrightarrow{} K_{t_{i+1}}
$$ 
 into  several simplex-wise inclusions in $\tilde{F}(\delta)$: 
$$  K_{t_{i-1}}=K'_{s} \incto K'_{t_i-\delta} \incto K'_{t_i}\incto K'_{t_i +\delta} \incto \cdots \incto K'_{t_{i+(k-1)\delta}} =K_{t_i}.
% \xrightarrow{} K_{t_{i+1}}.
$$ 
In the above, $K'_{t_i-\delta}$ is a ``dummy'' complex equal to $K'_{s}$
(so $K'_{s} \incto K'_{t_i-\delta}$ is an identity map)
and $k$ is the number of simplices added from $K_{t_{i-1}}$ to $K_{t_i}$ in $\tilde{F}$.
In short, adding the dummy complexes makes sure that the closed bars of the simplex-wise filtration 
``induce''
% approximate 
the closed-open time intervals in $\tilde{F}$ in which a cycle is harmonic; see \Cref{sec:stability} for details.
We also have that the $k$ number of simplices are added based on
their simplex ordering described above.

    
Let $span_{\delta}(z)$ denote the span of $z$ in $\tilde{F}(\delta)$, and analogously define $bar_{\delta}(z)$, $b_{\delta}(z)$, and $d_{\delta}(z)$. We observe that addition of dummy complexes implies that for a cycle $z$, bar of $z$ in $\tilde{F}(\delta)$ must approximate the span of $z$ in $\tilde{F}$; see Fig.~\ref{fig:spanbar} for an example. In more detail we have the following.

\begin{figure}
    \centering
    \includegraphics[width=0.7\linewidth]{spanbar.pdf}
    \caption{Turning closed-open spans into closed bars}
    \label{fig:spanbar}
\end{figure}

 
\begin{lemma}\label{lemma:barsnapping}
Let $m$ be the size of $K$
and let $z\in Z_*(K)$ be a cycle. For $\delta>0$ sufficiently small,
$t_{b(z)} \leq t_{b_\delta(z)} \leq t_{b(z)} +m\delta $, and $t_{d(z)} \leq t_{d_\delta(z)} \leq t_{d(z)} +m\delta$. Moreover, we have $ t_{d(z)} -\delta \leq t_{d_\delta(z)-1} \leq t_{d(z)}+(m-1)\delta$.

\end{lemma}



\begin{proof}

Since $K_{t_i-\delta}= K_{t_{i-1}}$, we have $b_\delta (z) \geq b(z)$.  It is clear that $b_\delta(z) \leq b(z)+m\delta$. 
The inequalities for $d_\delta(z)$ can also be easily checked. 

$ t_{d(z) -1} $ is the largest (real valued) index smaller than $t_d(z)$. The smallest that this index can be is the index of a dummy complex added before $K_{t_{d(z)}}$, and the largest it can be is the one to last index among complexes added after $K_{d(z)}$. 
\end{proof}



Therefore, the barcode of $\tilde{F}(\delta)$ can induce
% be taken as an approximation of 
a barcode for $\tilde{F}$.  Indeed, for any $\delta$ small enough, there is an obvious mapping $\Gamma: \tilde{F}(\delta) \xrightarrow{} \tilde{F}$. 
\begin{definition}
Let $F$ be a simplex-wise filtration (as in \Cref{sec:persistence}) and
let $\obarc(F)$ be the \emph{closed-open barcode of $\har{(F)}$}, i.e., $[t_b,t_{d}) \in \obarc(F)$ iff $[t_b,t_{d-1}] \in \hbarc(F)$.  
 Moreover, let $\tilde{F}$ and $\tilde{F}(\delta)$ be as defined in this section.
We define the \emph{harmonic chain barcode of $\tilde{F}$} as the set $\obarc(\tilde{F}) := \Gamma(\obarc(\tilde{F}(\delta)))$ for $\delta$ small enough.
\end{definition}

Note that $\obarc(\tilde{F})$ is the limit of $\hbarc(\tilde{F}(\delta))$ as $\delta$ approaches 0. 

In the rest of this section, we prove stability with respect to simplex-wise filtrations and at the end we show the stability for sublevel set harmonic barcode.




\subsection{$\mathbb{U}$-indexed Persistence Modules}\label{subsec:extension}


Let $\Rspace^{op}$ denote the poset $\Rspace$ with opposite ordering. Let $\mathbb{U}$ be the subset of $\Rspace^{op}\times \Rspace$ consisting of pairs $(a,b), a\leq b$. Note that $(a,b) < (c,d)$ if and only if $c<a, b < d$. Therefore, $\mathbb{U}$ can be considered as consisting of real-valued intervals with the order being inclusions between the intervals. A \emph{$\mathbb{U}$-indexed persistence module} $M$ is a functor from the poset $\mathbb{U}$ to the category of vector spaces. We have for each $(a,b)$, a vector space $M_{(a,b)}$, and for each relation $(a,b) < (c,d)$, a linear map\footnote{For simplicity we omit the subscripts for maps as long as this does not cause confusions.} $f:M_{(a,b)} \xrightarrow{} M_{(c,d)}$.

We now define the \emph{lift} (or \emph{extension}) of the harmonic zigzag module $\har(\filt)$, denoted $E= E(\har(\filt))$. We have $$ E(a,b) = \{ z \in Z_*(K)\mid b(z) \leq b ,  d(z)  \geq a) \}.$$ That is, $E(a,b)$ contains the chains that are born at or before $b$ and die at or after $a$ (see \Cref{dfn:bd-span-bar2}). Observe that if $z$ and $z'$ are born at or before $t$, then $z+z'$ is born at or before $t$, and if $z$ and $z'$ die at or after $t$, then $z+z'$ dies at or after $t$. To see the last statement, let $d$ be the minimum  time that $z$ and $z'$ die, and assume $z_1+z_2$ dies at $t_i < d$. Then $z+z' \notin K_{t_{i+1}} \subset K_{t_i}$. Therefore, one of $z$ and $z'$ must also not be in the image, which is a contradiction. From these observations we deduce the following.   


\begin{lemma}
    For any $(a,b) \in \mathbb{U}$, $E(a,b)$ is a vector space. 
\end{lemma}

\begin{lemma}\label{l:directsum}
    The extension operator $E$ respects direct sums. 
\end{lemma}


If $(a,b) < (c,d) $, then $[a,b] \subset [c,d]$ and we have an inclusion $E(a,b) \subset E(c,d)$. These are maps of the extension $\mathbb{U}$-module. Note that \textit{each} cycle of $K$ is in some $E(a,b)$ 
where $a$ could be possibly equal to $b$.
% perhaps on the diagonal.


\subsection{Blocks and Block-Decomposable Modules}


An \emph{interval} in a $\mathbb{U}$-Module $M$ is a subset ${\cal{J}} \subset \mathbb{U}$ such that 
\begin{enumerate}
    \item ${\cal{J}}$ is non-empty,
    \item if $p,q \in {\cal{J}}$ and $p \leq r \leq q$, then $r \in {\cal{J}},$ and,
    \item for any $p,q \in {\cal{J}}$, there is a sequence $p = r_0, r_1, \ldots, r_l = q$ of elements of ${\cal{J}}$ such that for all $0 \leq j \leq l-1$, $r_j$ and $r_{j+1}$ are comparable (connectivity). 
\end{enumerate}

A multiset of intervals in $\mathbb{U}$ is called a \emph{barcode}.

For ${\cal{J}}$ an interval in $\mathbb{U}$, the \emph{interval module $I^{{\cal{J}}}$} is a $\mathbb{U}$-indexed module such that $I^{{\cal{J}}}_p = \Rspace$ if $p \in {\cal{J}}$, and $I_a^{{\cal{J}}}=0$ otherwise. The maps satisfy $f:p\xrightarrow{} q = id_{\Rspace}$ if $p \leq q \in {\cal{J}}$ and $0$ otherwise.


The intervals for a $\mathbb{U}$-Module are indecomposable modules analogous to the interval modules in an $\Rspace$-indexed persistence module~\cite[Prop. 2.2]{botnan2018algebraic}. However, contrary to the 1D persistence, a general $\mathbb{U}$-indexed module cannot be  decomposed into interval modules. Note that the harmonic zigzag module is decompasable into interval modules. This implies that the module $E$ is decomposable into special forms of intervals called blocks.

For a closed interval $[a,b] \subset \Rspace^+$ the \emph{closed block} ${\mathcal{J}_{BL}} = [a,b]_{BL}$ is an interval in $\mathbb{U}$ defined by $$ [a,b]_{BL} = \{ (x,y) \in \mathbb{U}\mid  x \leq b,  y \geq a \}.$$ See~\cref{fig:block-extension}. 
A multiset of blocks in $\mathbb{U}$ is called a \emph{block barcode}. 




For any block ${\mathcal{J}_{BL}}$, the interval module $I^{{\mathcal{J}_{BL}}}$ is called a \emph{block module}. 
A $\mathbb{U}$-indexed module is \emph{block decomposable} if it can be written as a direct sum of block modules. 

\begin{figure}
    \centering
    \includegraphics[width=0.5\linewidth]{closed-block.pdf}
    \caption{The extension of a closed interval is a closed block.}
    \label{fig:block-extension}
\end{figure}


It is not difficult to check that if $[a,b]$ is a closed interval for a zigzag module, then $E(a,b)$ is a closed block module. Since the direct sum is preserved by $E$, we have  the following.


\begin{lemma}
    The $\mathbb{U}$-indexed module $E(\har(\filt))$ is block decomposable. 
\end{lemma}


\subsection{Interleaving of $\mathbb{U}$-Indexed Modules}







%In a $\mathbb{U}$-module a closed block is a block of the form 
To define $\eps$-interleaving between $\mathbb{U}$-indexed modules, we need to first define a \emph{shift operator} $S^\eps$. For a $\mathbb{U}$-indexed module $M$, we have a shifted module $S^\eps(M)$ such that 
$S^\eps(M)_{(a,b)} =M_{(a-\eps, b+\eps)}$, and for $f:M_{(c,d)} \xrightarrow{} M_{(a,b)}$, $S^\eps(f) = f: M_{(c-\eps,d+\eps)}\xrightarrow{} M_{(a-\eps,b+\eps) )}$.

% To define $\eps$-interleaving between $\mathbb{U}$-indexed modules, we need to first define a \emph{shift operator} $S^\eps$. For a $\mathbb{U}$-indexed module $M$, we have a shifted module $S^\eps(M)$ such that $S^\eps(M)_{(a,b)} = (a-\eps, b+\eps)$, and $f:M_{(c,d)} \xrightarrow{} M_{(a,b)}$ in $M$ maps to  $S^\eps(f): M_{(c-\eps,d+\eps)}\xrightarrow{} M_{(a-\eps,b+\eps) )}$.

Given $\eps \geq 0$, two $\mathbb{U}$-indexed modules $M ,N$ are \emph{$\eps$-interleaved} if there exist a pair of morphisms $\Phi: M \xrightarrow{} S^\eps (N)$ and $\Psi: N \xrightarrow{} S^\eps (M)$ such that for each $(a,b) \in \mathbb{U}$

$$ \Psi_{(a-\eps, b+\eps)}  \circ \Phi_{(a,b)} = f:M_{(a,b)} \xrightarrow{} M_{(a-2\eps, b+2\eps)}  $$
and
$$ \Phi_{(a-\eps, b+\eps)}  \circ \Psi_{(a,b)} = g:N_{(a,b)} \xrightarrow{} N_{(a-2\eps, b+2\eps)},$$
where  $g$ in the above denotes  maps of $N$.

The \emph{interleaving-distance} between two $\mathbb{U}$-indexed modules $M$ and $N$ is defined as $$ d_I(M,N) = \inf \{\eps\mid \text{there is an $\eps$-interleaving between $M$ and $N$} \}.$$

%\begin{proposition}

%Let $\har_*(F)$ and $\har_*(G)$ be two harmonic zigzag modules. Assume there exists a harmonic-preserving $\eps$-interleaving between $\har_*(F)$ and $\har_*(G)$. Then there is an $\eps$-interleaving between $E(\har_*(F))$ and $E(\har_*(G))$.
%\end{proposition}



\subsection{Stability of Harmonic Chain Barcode}\label{subsec:stability}


Let $F$ be a filtration. Define $\rho(F) = \max \{ t_{d(z)} - t_{d(z)-1} \mid z\in Z_*(K)\}$, that is, $\rho(F)$ the maximum distortion caused by turning spans into closed bars.


\begin{theorem}

Let $F$ and $G$ be simplex-wise increasing filtrations where the maps of $F$ and $G$ are inclusions. Assume that the interleaving distance at chain level between $F$ and $G$ is obtained via inclusions. 
Then $$ d_B( {\hbarc(F)}, {\hbarc(G)} ) \leq d_{I}(F,G) +\max \{\rho(F), \rho(G) \}.$$
\end{theorem}


\begin{proof}
    
    Set $M=E(\har(F))$ and $N=E(\har(G))$.  Let $\eps = d_{I}(F,G)$. We define $\Phi_{(a,b)}: M_{(a,b)} \xrightarrow{} N_{(a-\eps-\rho(F), b+\eps+\rho(F))}$ to be the inclusion. Note that $\Phi$ is well-defined since all the cycles exist in the extension. %, so we do not need to show that the image consists of only harmonic cycles. 
    %This is well-defined since each element of $M_{(a,b)}$ is a cycle $z$ whose span in $F$ is included in $[a,b+\rho(F)]$, and by~\cref{lemma:box-lemma}, the span of $z$ in $G$ is included in $[a-\eps-\rho(F), b+\eps+\rho(F)]$. 
    The morphism $\Psi = \{ \Psi_{(a,b)}\mid (a,b)\in \mathbb{U} \}$ is defined analogously. Since all the maps are inclusions, the morphisms $\Phi$ and $\Psi$ define an $\eps$-interleaving of $M$ and $N$. It follows that $d_I (M,N) \leq \eps+\max \{\rho(F), \rho(G) \}$.
    From the isometry theorem for block decomposable $\mathbb{U}$-indexed modules~\cite{bjerkevik2021stability,botnan2018algebraic}, we deduce that $d_B( {\mathcal{D}} (M), {\mathcal{D}} (N)) = d_I(M,N) \leq \eps +\max \{\rho(F), \rho(G) \}$, where ${\mathcal{D}}(M)$ and ${\mathcal{D}}(N)$ are the block barcodes of $M$ and $N$ respectively. Any matching of block barcodes induces a matching of the parts of the blocks on the diagonal which form the barcodes of the 1D filtrations $\har(F)$ and $\har(G)$ and vice versa (\cite[Proposition 4.2]{botnan2018algebraic}). It follows that $d_B (\hbarc(F), \hbarc(G)) = d_B( {\mathcal{D}} (M), {\mathcal{D}} (N))  \leq \eps + \max \{\rho(F), \rho(G)\}$. 
\end{proof}



\begin{corollary}
Let $\hat{f}, \hat{g} : |K| \xrightarrow{} \Rspace$ be simplex-wise linear functions, and let $\tilde{F}$ and $\tilde{G}$ be the sublevel set filtrations of $\hat{f}$ and $\hat{g}$ respectively. Then, $$d_{B}({\obarc(\tilde{F})}, {\obarc(\tilde{G}})) \leq || \hat{f}-\hat{g}||_{\infty}.$$
\end{corollary}

\begin{proof}
    %Note that $||f-g||_\infty = \max \{ |f(v) - g(w)|, v,w \in V(K) \}$. For $K_t$ in $F$ and $K'_s$ in $G$ we have $|t-s| \leq ||f-g||_\infty$. 

    %We show that for small enough $\delta$,  $d_{B}({\obarc(\tilde{F}(\delta))}, {\obarc(\tilde{G}(\delta)}) \leq || f-g||_{\infty}$ from which the corollary follows.
    
    Let $m$ be the number of simplices in $K$. Let $\eps = ||\hat{f}-\hat{g}||_\infty$. There is an $\eps$-interleaving between the filtration $\tilde{F}$ and $\tilde{G}$ defined by inclusion of sublevel sets~\cite{Cohen-SteinerEdelsbrunnerHarer2005,ChazalCohenSteinerGlisse2009}. We can choose $\delta$ such that $\eps \gg m\delta$.
    %then $\eps$-interleaving between sublevel set filtrations $\tilde{F}$ and $\tilde{G}$ implies an $\eps$-interleaving between $\tilde{F}(\delta)$ and $\tilde{G}(\delta)$.
    Then there is an $(\eps+m\delta)$-interleaving between $\tilde{F}(\delta)$ and $\tilde{G}(\delta)$, since $K_{t} \subset K'_{t+\eps+m\delta}$ and $K'_{t} \subset K_{t+\eps+m\delta}$, where $K_t$ and $K'_t$ are the complexes of $\tilde{F}(\delta)$ and $\tilde{G}(\delta)$ respectively. Moreover, by~\cref{lemma:barsnapping}, $\max\{\rho(F), \rho(G)\} \leq m\delta$. Therefore, for all $\delta$ such that $m\delta \ll \eps$, $d_B ( {\obarc(\tilde{F}(\delta))}, {\obarc(\tilde{G}(\delta))} )  \leq \eps+m\delta = ||f-g||_\infty+m\delta$. Since $\delta$ approaches 0 the corollary follows. 
    %For small enough $\delta>0$, the complexes
    %$K_{t}$ in $F(\delta)$ and $K'_s$ in $G(\delta)$ satisfy $|t-s| \geq ||f-g||+$
\end{proof}





%\subsection{Stability Example}

%We consider the same example we presented for instability, see~\cref{fig:stableexample}. 
%For the left filtration, at time $t_4$, the cycle $z_1 = t_4+t_2-t_1$ is created and represents the top bar. 
%At time $t_5$ the dimension of the harmonic chains increases, and one can see that the cycle $z_2 = t_2-t_1+t_3-t_5$ is harmonic at $t_5$ and survives up to $t_9$.
%Among the cycles that exist at $t_5$, a 1-dimensional subspace survives $t_8$, this space is generated by $z_2$.
%\Bei{?}
%For cycles that live at time $t_5$, after adding $t_8$ 
%Among edges inserted until $t_5$, no cycle survives adding $t_8$ other than multiples of $z_2$. 
%At $t_8$, the harmonic cycle generate 
%\Bei{Please add a sentence here.}
%At $t_7$, all edges are present, therefore there has to be a cycle that survives until $t_{10}$. 
%We see that $z_3 = \partial(t_{10}) - 1/2 \partial (t_8 - t_9 )$ has this property and serves as the harmonic representative.

%For the right filtration, at time $t_4$, the cycle $\bar{z}_1 = -t_1+t_2+t_4-t_3$ is created and survives up to $t_9$. 
%At time $t_5$, $\bar{z}_2 = t_5+t_4-t_3$ is created and  survives up to $t_8$. 
%Finally, at time $t_7$, $\bar{z}_3=z_3$ is created and survives up to $t_{10}$.

%We observe that the first bar on top on the left can be matched to the second bar on the right, and the second bar in the left can be 
%matched to the first bar on the right. The bottom bar is unchanged in both barcodes.

%\begin{figure}
%    \centering
%    \includegraphics{stable-example.pdf}
%    \caption{The canonical barcodes of harmonic chains (in blue) are stable.}
%    \label{fig:stableexample}
%\end{figure}
% END AUTHOR FILE sec-stability.tex
%\input{sec-conclusion}
% BEGIN AUTHOR FILE sec-acknowledgments.tex
\section*{Acknowledgment}
T.\ Hou is supported by NSF fund CCF 2439255. B.\ Wang is supported by DOE DE-SC0021015, NSF IIS-2145499 and NSF DMS-2301361. 
% END AUTHOR FILE sec-acknowledgments.tex
%------------------------------
% BEGIN AUTHOR FILE Arxiv-fullversion-2025.bbl
\begin{thebibliography}{10}

\bibitem{bjerkevik2021stability}
H{\aa}vard Bakke~Bjerkevik.
\newblock On the stability of interval decomposable persistence modules.
\newblock {\em Discrete \& Computational Geometry}, 66(1):92--121, 2021.

\bibitem{BasuCox2022}
Saugata Basu and Nathanael Cox.
\newblock Harmonic persistent homology.
\newblock In {\em 2021 IEEE 62nd Annual Symposium on Foundations of Computer
  Science (FOCS)}, pages 1112--1123. IEEE, 2022.

\bibitem{bauer2014induced}
Ulrich Bauer and Michael Lesnick.
\newblock Induced matchings of barcodes and the algebraic stability of
  persistence.
\newblock In {\em Proceedings of the thirtieth annual symposium on
  Computational geometry}, pages 355--364, 2014.

\bibitem{BendichPaul2013Haro}
Paul Bendich, Herbert Edelsbrunner, Dmitriy Morozov, and Amit Patel.
\newblock Homology and robustness of level and interlevel sets.
\newblock {\em Homology, homotopy, and applications}, 15(1):51--72, 2013.

\bibitem{Borradaile2020}
Glencora Borradaile, William Maxwell, and Amir Nayyeri.
\newblock Minimum bounded chains and minimum homologous chains in embedded
  simplicial complexes.
\newblock In {\em Proceedings of the 36th International Symposium on
  Computational Geometry (SoCG 2020)}. Schloss Dagstuhl-Leibniz-Zentrum f{\"u}r
  Informatik, 2020.

\bibitem{botnan2018algebraic}
Magnus Botnan and Michael Lesnick.
\newblock Algebraic stability of zigzag persistence modules.
\newblock {\em Algebraic \& geometric topology}, 18(6):3133--3204, 2018.

\bibitem{Oleksiy2010}
Oleksiy Busaryev, Tamal~K. Dey, and Yusu Wang.
\newblock Tracking a generator by persistence.
\newblock In My~T. Thai and Sartaj Sahni, editors, {\em Computing and
  Combinatorics}, pages 278--287, 2010.

\bibitem{CarlssonSilva2010}
Gunnar Carlsson and Vin de~Silva.
\newblock Zigzag persistence.
\newblock {\em Foundations of Computational Mathematics}, 10(4):367--405, 2010.

\bibitem{CarlssonSilvaMorozov2009}
Gunnar Carlsson, Vin de~Silva, and Dmitriy Morozov.
\newblock Zigzag persistent homology and real-valued functions.
\newblock In {\em Proceedings of the 25th Annual Symposium on Computational
  Geometry}, pages 247--256, 2009.

\bibitem{CarlssonZomorodianCollins2004}
Gunnar Carlsson, Afra~J. Zomorodian, Anne Collins, and Leonidas~J. Guibas.
\newblock Persistence barcodes for shapes.
\newblock {\em Proceedings Eurographs/{ACM} {SIGGRAPH} Symposium on Geometry
  Processing}, pages 124--135, 2004.

\bibitem{chambers2019minimum}
Erin~W. Chambers, Jeff Erickson, Kyle Fox, and Amir Nayyeri.
\newblock Minimum cuts in surface graphs.
\newblock {\em SIAM Journal on Computing}, 52(1):156--195, 2023.
\newblock \href {http://arxiv.org/abs/https://doi.org/10.1137/19M1291820}
  {\path{arXiv:https://doi.org/10.1137/19M1291820}}, \href
  {https://doi.org/10.1137/19M1291820} {\path{doi:10.1137/19M1291820}}.

\bibitem{Chambers2009}
Erin~W. Chambers, Jeff Erickson, and Amir Nayyeri.
\newblock Minimum cuts and shortest homologous cycles.
\newblock In {\em Proceedings of the 25th International Symposium on
  Computational Geometry ({SoCG} 2009)}. {ACM} Press, 2009.
\newblock \href {https://doi.org/10.1145/1542362.1542426}
  {\path{doi:10.1145/1542362.1542426}}.

\bibitem{ChambersParsaSchreiber2022}
Erin~Wolf Chambers, Salman Parsa, and Hannah Schreiber.
\newblock On complexity of computing bottleneck and lexicographic optimal
  cycles in a homology class.
\newblock In {\em Proceedings of the 38th International Symposium on
  Computational Geometry (SoCG)}. Schloss Dagstuhl-Leibniz-Zentrum f{\"u}r
  Informatik, 2022.

\bibitem{ChazalCohenSteinerGlisse2009}
Fr{\'{e}}d{\'{e}}ric Chazal, David Cohen-Steiner, Marc Glisse, Leonidas~J.
  Guibas, and Steve~Y. Oudot.
\newblock Proximity of persistence modules and their diagrams.
\newblock In {\em 25th Annual Symposium on Computational Geometry (SoCG 2009)},
  pages 237--246, New York, NY, USA, 2009. Association for Computing Machinery.
\newblock \href {https://doi.org/https://doi.org/10.1145/1542362.1542407}
  {\path{doi:https://doi.org/10.1145/1542362.1542407}}.

\bibitem{chazal2016structure}
Fr{\'e}d{\'e}ric Chazal, Vin De~Silva, Marc Glisse, and Steve Oudot.
\newblock {\em The Structure and Stability of Persistence Modules}, volume~10.
\newblock Springer, 2016.

\bibitem{Chen2010}
Chao Chen and Daniel Freedman.
\newblock Measuring and computing natural generators for homology groups.
\newblock {\em Computational Geometry}, 43(2):169--181, 2010.
\newblock Special Issue on the 24th European Workshop on Computational Geometry
  (EuroCG'08).
\newblock \href {https://doi.org/10.1016/j.comgeo.2009.06.004}
  {\path{doi:10.1016/j.comgeo.2009.06.004}}.

\bibitem{ChenFreedman2011}
Chao Chen and Daniel Freedman.
\newblock Hardness results for homology localization.
\newblock {\em Discrete \& Computational Geometry}, 45(3):425--448, 2011.

\bibitem{ChenQiuWang2022}
Jiahui Chen, Yuchi Qiu, Rui Wang, and Guo-Wei Wei.
\newblock Persistent {L}aplacian projected {O}micron {BA. 4} and {BA. 5} to
  become new dominating variants.
\newblock {\em Computers in Biology and Medicine}, 151:106262, 2022.

\bibitem{Cohen-SteinerEdelsbrunnerHarer2005}
David Cohen-Steiner, Herbert Edelsbrunner, and John Harer.
\newblock Stability of persistence diagrams.
\newblock In {\em Proceedings of the 21st Annual Symposium on Computational
  Geometry}, pages 263--271, 2005.

\bibitem{cohen2006vines}
David Cohen-Steiner, Herbert Edelsbrunner, and Dmitriy Morozov.
\newblock Vines and vineyards by updating persistence in linear time.
\newblock In {\em Proceedings of the Twenty-Second Annual Symposium on
  Computational Geometry}, pages 119--126, 2006.

\bibitem{DeGregorioGuerraScaramuccia2021}
Alessandro De~Gregorio, Marco Guerra, Sara Scaramuccia, and Francesco
  Vaccarino.
\newblock Parallel decomposition of persistence modules through interval bases.
\newblock {\em arXiv preprint arXiv:2106.11884}, 2021.

\bibitem{DeSilvaMorozovVejdemoJohansson2011}
Vin De~Silva, Dmitriy Morozov, and Mikael Vejdemo-Johansson.
\newblock Persistent cohomology and circular coordinates.
\newblock {\em Discrete \& Computational Geometry}, 45(4):737--759, 2011.

\bibitem{DeyHiraniKrishnamoorthy2010}
Tamal~K. Dey, Anil~N. Hirani, and Bala Krishnamoorthy.
\newblock Optimal homologous cycles, total unimodularity, and linear
  programming.
\newblock In {\em Proceedings of the 42nd ACM Symposium on Theory of
  Computing}, pages 221--230, 2010.
\newblock \href {https://doi.org/10.1145/1806689.1806721}
  {\path{doi:10.1145/1806689.1806721}}.

\bibitem{dey2021computing}
Tamal~K. Dey and Tao Hou.
\newblock Computing zigzag persistence on graphs in near-linear time.
\newblock In {\em 37th International Symposium on Computational Geometry},
  2021.

\bibitem{DeyHouMandal2020}
Tamal~K Dey, Tao Hou, and Sayan Mandal.
\newblock Computing minimal persistent cycles: Polynomial and hard cases.
\newblock In {\em Proceedings of the 14th Annual {ACM-SIAM} Symposium on
  Discrete Algorithms}, pages 2587--2606. SIAM, 2020.

\bibitem{zzrep}
Tamal~K. Dey, Tao Hou, and Dmitriy Morozov.
\newblock A fast algorithm for computing zigzag representatives.
\newblock {\em ACM-SIAM Symposium on Discrete Algorithms (SODA) 2025, {\sf to
  appear}}, 2025.

\bibitem{DeyWang2022}
Tamal~Krishna Dey and Yusu Wang.
\newblock {\em Computational Topology for Data Analysis}.
\newblock Cambridge University Press, 2022.

\bibitem{Eckmann1944}
Beno Eckmann.
\newblock Harmonische {F}unktionen und {R}andwertaufgaben in einem {K}omplex.
\newblock {\em Commentarii Mathematici Helvetici}, 17(1):240--255, 1944.

\bibitem{EdelsbrunnerLetscherZomorodian2002}
Edelsbrunner, Letscher, and Zomorodian.
\newblock Topological persistence and simplification.
\newblock {\em Discrete \& Computational Geometry}, 28:511--533, 2002.

\bibitem{EdelsbrunnerHarer2010}
Herbert Edelsbrunner and John Harer.
\newblock {\em Computational Topology: An Introduction}.
\newblock Applied Mathematics. American Mathematical Society, 2010.

\bibitem{Ghrist2008}
Robert Ghrist.
\newblock Barcodes: The persistent topology of data.
\newblock {\em Bullentin of the American Mathematical Society}, 45:61--75,
  2008.

\bibitem{Grochow2018}
Joshua~A. Grochow and Jamie Tucker-Foltz.
\newblock Computational topology and the unique games conjecture.
\newblock In {\em Proceedings of 34th International Symposium on Computational
  Geometry (SoCG 2018)}. Schloss Dagstuhl-Leibniz-Zentrum fuer Informatik,
  2018.

\bibitem{GuglielmiSavostianovTudisco2023}
Nicola Guglielmi, Anton Savostianov, and Francesco Tudisco.
\newblock Quantifying the structural stability of simplicial homology.
\newblock {\em arXiv preprint arXiv:2301.03627}, 2023.

\bibitem{GulenMemoliWan2024}
Aziz~Burak G{\"u}len, Facundo M{\'e}moli, and Zhengchao Wan.
\newblock Orthogonal {M}\"obius inversion and {G}rassmannian persistence
  diagrams.
\newblock arXiv preprint arXiv:2311.06870, 2024.

\bibitem{GurnariGuzman-SaenzUtro2023}
Davide Gurnari, Aldo Guzm{\'a}n-S{\'a}enz, Filippo Utro, Aritra Bose, Saugata
  Basu, and Laxmi Parida.
\newblock Probing omics data via harmonic persistent homology.
\newblock arXiv preprint arXiv:2311.06357, 2023.

\bibitem{hatcher2002algebraic}
Allen Hatcher.
\newblock {\em Algebraic Topology}.
\newblock Cambridge University Press, 2002.

\bibitem{hiraoka2016hierarchical}
Yasuaki Hiraoka, Takenobu Nakamura, Akihiko Hirata, Emerson~G Escolar, Kaname
  Matsue, and Yasumasa Nishiura.
\newblock Hierarchical structures of amorphous solids characterized by
  persistent homology.
\newblock {\em Proceedings of the National Academy of Sciences},
  113(26):7035--7040, 2016.

\bibitem{HorakJost2013}
Danijela Horak and J{\"u}rgen Jost.
\newblock Spectra of combinatorial {L}aplace operators on simplicial complexes.
\newblock {\em Advances in Mathematics}, 244:303--336, 2013.

\bibitem{KerosSubr2023}
Alexandros Keros and Kartic Subr.
\newblock Spectral coarsening with {H}odge {L}aplacians.
\newblock In {\em ACM SIGGRAPH 2023 Conference Proceedings}, pages 1--11, 2023.

\bibitem{Kirchhoff1847}
G.~Kirchhoff.
\newblock Ueber die {A}ufl{\"o}sung der {G}leichungen, auf welche man bei der
  {U}ntersuchung der linearen {V}ertheilung galvanischer {S}tr{\"o}me
  gef{\"u}hrt wird.
\newblock {\em Annalen der Physik}, 148(12):497--508, 1847.
\newblock \href {https://doi.org/10.1002/andp.18471481202}
  {\path{doi:10.1002/andp.18471481202}}.

\bibitem{LeeChungChoi2019}
Hyekyoung Lee, Moo~K Chung, Hongyoon Choi, Hyejin Kang, Seunggyun Ha, Yu~Kyeong
  Kim, and Dong~Soo Lee.
\newblock Harmonic holes as the submodules of brain network and network
  dissimilarity.
\newblock In {\em Computational Topology in Image Context: 7th International
  Workshop, CTIC 2019, M{\'a}laga, Spain, January 24-25, 2019, Proceedings 7},
  pages 110--122. Springer, 2019.

\bibitem{LiDuncanTopp2017}
Mao Li, Keith Duncan, Christopher~N. Topp, and Daniel~H. Chitwood.
\newblock Persistent homology and the branching topologies of plants.
\newblock {\em American Journal of Botany}, 104(3):349--353, 2017.

\bibitem{Lieutier}
Andre Lieutier.
\newblock Persistent harmonic forms.
\newblock
  \url{https://project.inria.fr/gudhi/files/2014/10/Persistent-Harmonic-Forms.pdf}.

\bibitem{LiuLiWu2023}
Jian Liu, Jingyan Li, and Jie Wu.
\newblock The algebraic stability for persistent {L}aplacians.
\newblock {\em arXiv preprint arXiv:2302.03902}, 2023.

\bibitem{LuoHenselman-Petrusek2023}
Jiajie Luo and Gregory Henselman-Petrusek.
\newblock Interval decomposition for persistence modules freely generated over
  {PIDs}.
\newblock arXiv preprint arXiv:2310.07971, 2023.

\bibitem{maria2015zigzag}
Cl{\'{e}}ment Maria and Steve~Y. Oudot.
\newblock Zigzag persistence via reflections and transpositions.
\newblock In Piotr Indyk, editor, {\em Proceedings of the Twenty-Sixth Annual
  {ACM-SIAM} Symposium on Discrete Algorithms, {SODA} 2015, San Diego, CA, USA,
  January 4-6, 2015}, pages 181--199. {SIAM}, 2015.

\bibitem{MemoliWanWang2022}
Facundo M{\'e}moli, Zhengchao Wan, and Yusu Wang.
\newblock Persistent {L}aplacians: Properties, algorithms and implications.
\newblock {\em SIAM Journal on Mathematics of Data Science}, 4(2):858--884,
  2022.

\bibitem{Merris1994}
Russell Merris.
\newblock {L}aplacian matrices of graphs: a survey.
\newblock {\em Linear Algebra and its Applications}, 197-198:143--176, 1994.
\newblock \href {https://doi.org/10.1016/0024-3795(94)90486-3}
  {\path{doi:10.1016/0024-3795(94)90486-3}}.

\bibitem{MikePerea2019}
Joshua Mike and Jose Perea.
\newblock Multiscale geometric data analysis via {L}aplacian eigenvector
  cascading.
\newblock In {\em 2019 18th IEEE International Conference On Machine Learning
  And Applications (ICMLA)}, pages 1091--1098. IEEE, 2019.

\bibitem{MoharAlaviChartrand1991}
Bojan Mohar, Y~Alavi, G~Chartrand, and OR~Oellermann.
\newblock The {L}aplacian spectrum of graphs.
\newblock {\em Graph theory, combinatorics, and applications}, 2(871-898):12,
  1991.

\bibitem{Obayashi2018}
Ippei Obayashi.
\newblock Volume-optimal cycle: Tightest representative cycle of a generator in
  persistent homology.
\newblock {\em SIAM Journal on Applied Algebra and Geometry}, 2(4):508--534,
  2018.
\newblock \href {https://doi.org/10.1137/17M1159439}
  {\path{doi:10.1137/17M1159439}}.

\bibitem{WangNguyenWei2020}
Rui Wang, Duc~Duy Nguyen, and Guo-Wei Wei.
\newblock Persistent spectral graph.
\newblock {\em International journal for numerical methods in biomedical
  engineering}, 36(9):e3376, 2020.

\bibitem{Zomorodian2005}
Afra~J Zomorodian.
\newblock {\em Topology for computing}, volume~16.
\newblock Cambridge university press, 2005.

\end{thebibliography}
% END AUTHOR FILE Arxiv-fullversion-2025.bbl
%------------------------------
\appendix  
% BEGIN AUTHOR FILE sec-homology.tex
\section{Homology and Cohomology}
\label{sec:homology}

Let $K$ be a simplicial complex and we give the standard orientation to simplices in $K$. We write simplices as an ordered set of vertices, e.g.,~$\sigma=v_0v_1\cdots v_{p}$ for a $p$-dimensional simplex, for any integer $p \geq 0$.
Let $K_p$ denote the set of $p$-simplices of $K$ and $n_p=|K_p|$. 
The $p$-dimensional \emph{chain group} of $K$ with coefficients in $\Rspace$, denoted $C_p(K)$, is an $\Rspace$-vector space generated by $K_p$. 
By fixing an ordering of the set $K_p$, we can identify any $p$-chain $c\in C_p(K)$ with an ordered $n_p$-tuple with real entries (e.g.,~an element of $\Rspace^{n_p}$).
For each $p$, we fix an ordering for the $p$-simplices once and for all, and identify $C_p(K)$ and $\Rspace^{n_p}$. 
The standard basis of $\Rspace^{n_p}$ corresponds (under the identification) to the basis of $C_p(K)$ given by the simplices with standard orientation.

The $p$-dimensional boundary matrix $\partial_p: C_p(K) \rightarrow C_{p-1}(K)$ is defined on a simplex basis element by the formula
\[ \partial (v_0v_1 \cdots v_p)= \sum_{j=0}^{p} {(-1)^q (v_0\cdots v_{q-1}v_{q+1}\cdots v_p) }.\]
In the right hand side above, in the $q$-th term $v_q$ is dropped. The formula guarantees the crucial property of the boundary homomorphism: $\partial_{p} \partial_{p+1} = 0$. This simply means that the boundary of a simplex has no boundary. The sequence $C_p(K)$ together with the maps $\partial_p$ define the \textit{simplicial chain complex} of $K$ with real coefficients, denoted $C_{\bullet}(K)$. The group of $p$-dimensional \textit{cycles}, denoted $Z_p(K)$, is the kernel of $\partial_p$. The group of $p$-dimensional boundaries, denoted $B_p(K)$, is the image of $\partial_{p+1}$. The $p$-dimensional homology group of $K$, denoted $H_p(K)$, is the quotient group $Z_p(K)/B_p(K)$. 
As a set, this quotient is formally defined as $\{z+B_p(K) \mid  z\in Z_p(K)\}$. The operations are inherited from the chain group. In words, the homology group is obtained from $Z_p(K)$ by setting any two cycles which differ by a boundary to be equal. All these groups are $\Rspace$-vector spaces.

Consider the space $\Rspace^{n_p}$. The cycle group $Z_p(K) \subset C_p(K) = \Rspace^{n_p}$ is a subspace, that is, a hyperplane passing through the origin. Similarly, the boundary group $B_p(K)$ is a subspace, and is included in $Z_p(K)$. The homology group is the set of parallels of $B_p(K)$ inside $Z_p(K)$. Each such parallel hyperplane differs from $B_p(K)$ by a translation given by some cycle $z$. These parallel hyperplanes partition $Z_p(K)$. The homology group is then isomorphic to the subspace perpendicular to $B_p(K)$ inside $Z_p(K)$. 
The dimension of the $p$-dimensional homology group is called the $p$-th $\Rspace$-Betti number, denoted $\beta_p(K)$. 

Simplicial \textit{cohomology} with coefficients in $\Rspace$ is usually  defined by the process of dualizing. This means that we replace an $p$-chain by a linear functional $C_p(K) \rightarrow \Rspace$, called a $p$-dimensional \textit{cochain}. 
The set of all such linear functionals is a vector space isomorphic to $C_p(K)$, called the cochain group, denoted $C^p(K)$, which is a dual vector space of $C_p(K)$.  
For the purposes of defining harmonic chains, we must fix an isomorphism. 
We take the isomorphism that sends each standard basis element of $\Rspace^{n_p}$, corresponding to $\sigma$, to a functional that assigns 1 to $\sigma$ and 0 to other basis elements, denoted $\hat{\sigma} \in C^p(K)$. Any cochain $\gamma \in C^p(K)$ can be written as a linear combination of the $\hat{\sigma}$. Therefore, it is also a vector in $\Rspace^{n_p}$. The fixed isomorphism allows us to identify $C^p$ with $\Rspace^{n_p}$, and hence to $C_p$. Therefore, any vector in $\Rspace^{n_p}$ is at the same time a chain and a cochain.

The \textit{coboundary matrix} $\delta^p: C^p(K)\rightarrow C^{p+1}(K)$, in its matrix representation, is the transpose of $\partial_{p+1}$, $\delta^p = \partial_{p+1}^\top$. It follows that $\delta_{p+1}\delta_p=0$, and we can form a \textit{cochain complex} $C^{\bullet}(K)$. The group of \textit{$p$-cocycles}, denoted $Z^p$ is the kernel of $\delta^p$. The group of \textit{$p$-coboundary} is the image of $\delta_{p-1}$. The $p$-dimensional \textit{cohomology group}, denoted $H^p(K)$, is defined as $H^p(K) = Z^p(K)/B^p(K)$. All of these groups are again vector subspaces of $C^p(K)$ and thus of $\Rspace^{n_p}$.
It is a standard fact that homology and cohomology groups with real coefficients are isomorphic.
% END AUTHOR FILE sec-homology.tex
% BEGIN AUTHOR FILE sec-rep-justify.tex
\newcommand{\img}{\mathrm{img}}
% \newcommand{\ker}{\mathrm{ker}}

\section{Justification for the Definition of Harmonic Representatives}\label{sec:rep-just}

Consider a zigzag module $\Mcal:V_0\lrarrowsp{g_0}V_1\lrarrowsp{g_1}\cdots\lrarrowsp{g_{k-1}}V_k$
% where each $V_i$ is a vector space and each $g_i$ is a linear map
% which could be forward, of the form $g_i:V_i\to V_{i+1}$,
% or backward, of the form $g_i:V_{i}\leftarrow V_{i+1}$.
% Furthermore, we assume that $\Mcal$
which is \emph{elementary} in the sense that each $g_i$
is either an isomorphism, an injection with corank $1$,
or a surjection with nullity $1$.
Representatives for intervals in $\barc(\Mcal)$
are defined as follows:

\begin{definition}[Zigzag representatives \cite{maria2015zigzag}]
\label{dfn:rep-cls-ori}
A \emph{representative} for $[b,d]\in\barc(\Mcal)$
is a sequence $\{v_i\in V_i\mid i\in[b,d]\}$
such that  for every $b\leq i<d$, 
either $g_i(v_i)=v_{i+1}$ or
$g_i(v_{i+1})=v_{i}$
based on the direction of $g_i$.
Furthermore, we have:
\begin{description}
% \label{dfn:rep}
    % \item \label{itm:dfn-rep-cls-birth}
    \item[Birth condition:] 
    If $g_{b-1}:V_{b-1}\to V_{b}$ is forward
    (thus being injective),
    $v_b\not\in\img(g_{b-1})$;
    % then $\cyc_b$ contains $\fsimp_{b-1}$;
    % then $[\cyc_b]$ is not in $\img(\linmap_{s-1})$;
    if $g_{b-1}:V_{b-1}\leftarrow V_{b}$ is backward
    (thus being  surjective),
    then $v_b$ is the non-zero element in $\ker(g_{b-1})$.
    
    % \item \label{itm:dfn-rep-cls-death}
    \item[Death condition:]
    If $d<k$ and
    $g_{d}:V_{d}\leftarrow V_{d+1}$ is backward
    (thus being injective),
    $v_d\not\in\img(g_{d})$;
    if $d<i$
    and $g_{d}:V_{d}\to V_{d+1}$ is forward
    (thus being surjective),
    then $v_d$ is the non-zero element in $\ker(g_{d})$.
\end{description}
% \smallskip
\end{definition}

\begin{proof}[Proof of \Cref{prop:single-rep}]
Let $[b,d]$ be an interval in $\hbarc_*(\filt)=\barc(\har_*(\filt))$
and let $\{z_i\in\har_*(K_i)\mid i\in[b,d]\}$
be a sequence of representatives for $[b,d]$
as in \Cref{dfn:rep-cls-ori}.
For each $i<d$,
we have that $z_i$
maps to $z_{i+1}$ 
or $z_{i+1}$
maps to $z_{i}$ 
by the inclusion,
and hence $z_i=z_{i+1}$.
Therefore, the representative
$\{z_i\in\har_*(K_i)\mid i\in[b,d]\}$ indeed
consists of a single cycle.
\end{proof}

To see that \Cref{dfn:rep} is an adaption of \Cref{dfn:rep-cls-ori}
to the module $\har_p(\filt)$, we only need to check the birth
and death conditions.
Let $[b,d]$ be an interval in $\hbarc_p(\filt)$
with a representative $z$.
For the birth condition, since inclusion maps have trivial kernels,
we have that $z$ is not in the image of the inclusion $\har_p(K_{b-1})\incto \har_p(K_b)$
by \Cref{dfn:rep-cls-ori}.
This means that 
$z\in\har_p(K_{b})\setminus\har_p(K_{b-1})$ 
which is as stated in \Cref{dfn:rep}.
The death condition can also be similarly verified.
% END AUTHOR FILE sec-rep-justify.tex
%\input{sec-barcode-zigzag-algo-impl}
%\input{sec-pf-prop-dim-chng}
% BEGIN AUTHOR FILE sec-harmonic-cycle-basis.tex
\section{Justification for $\Harmat^p$ Forming a Basis for $\har_p(K_i)$}\label{sec:harmat-basis}

We first have the following proposition which follows directly from definitions:
\begin{proposition}\label{prop:har-rep-basis}
For a $K_i$ in $\filt$,
let $\{[b_j,d_j]\mid j\in J\}$ 
be all the intervals in $\hbarc_p(\filt)$ containing $i$,
where each $[b_j,d_j]$ has a harmonic representative $z_j$.
Then $\{z_j\mid j\in J\}$ forms a basis for $\har_p(K_i)$.
\end{proposition}

Define a \emph{prefix} $\filt_i$ of $\filt$ as the filtration
% \begin{equation*}
% $
$\filt_i: 
% \emptyset
K_0 \incto K_1 \incto
\cdots \incto K_i.$
% $
% \end{equation*}
It is then not hard to see the following fact concerning \Cref{alg:frame}:
\begin{proposition}\label{prop:iter-unpair-to-i}
Before each iteration $i$
processing $K_{i}\incto K_{i+1}$,
$\{[b,i]\mid b\in U^p\}$
is the set of all intervals containing $i$
in $\hbarc_p(\filt_i)$
and each partial representative for $b$
maintained by the algorithm
is indeed a harmonic representative for $[b,i]\in \hbarc_p(\filt_i)$.
\end{proposition}

Combining \Cref{prop:har-rep-basis,prop:iter-unpair-to-i},
we have that columns in $\Harmat^p$ form a basis for $\har_p(K_i)$.
% END AUTHOR FILE sec-harmonic-cycle-basis.tex
